%% ============================================================================
%%  PRAMANIX — A Deterministic Neuro-Symbolic Execution Firewall
%%             for Autonomous AI Agents
%%
%%  A Technical Thesis: Engineering Design, Research Foundations,
%%  and Production Evidence
%%
%%  Author  : Viraj Jain
%%  Version : Pramanix 1.0.0
%%  Audit   : 2026-05-12
%%  License : AGPL-3.0-only
%% ============================================================================

\documentclass[12pt,a4paper,oneside]{report}

%% ── Encoding & fonts ────────────────────────────────────────────────────────
\usepackage[T1]{fontenc}
\usepackage[utf8]{inputenc}
\usepackage{lmodern}
\usepackage{microtype}

%% ── Page geometry ───────────────────────────────────────────────────────────
\usepackage[top=2.5cm, bottom=2.5cm, left=3.5cm, right=2.5cm]{geometry}

%% ── Mathematics ─────────────────────────────────────────────────────────────
\usepackage{amsmath, amssymb}

%% ── Tables ──────────────────────────────────────────────────────────────────
\usepackage{booktabs}
\usepackage{longtable}
\usepackage{array}
\usepackage{tabularx}

%% ── Code listings ───────────────────────────────────────────────────────────
\usepackage{listings}
\usepackage{xcolor}

\definecolor{codebg}{RGB}{248,248,248}
\definecolor{codeframe}{RGB}{210,210,210}
\definecolor{kwcolor}{RGB}{0,0,160}
\definecolor{cmtcolor}{RGB}{0,120,0}
\definecolor{strcolor}{RGB}{160,0,0}

\lstdefinestyle{pramanix}{
  backgroundcolor  = \color{codebg},
  basicstyle       = \small\ttfamily,
  breakatwhitespace= false,
  breaklines       = true,
  captionpos       = b,
  commentstyle     = \color{cmtcolor}\itshape,
  frame            = single,
  rulecolor        = \color{codeframe},
  keepspaces       = true,
  keywordstyle     = \color{kwcolor}\bfseries,
  language         = Python,
  numbers          = none,
  showspaces       = false,
  showstringspaces = false,
  stringstyle      = \color{strcolor},
  tabsize          = 4,
  upquote          = true,
}

\lstdefinestyle{plain}{
  backgroundcolor  = \color{codebg},
  basicstyle       = \small\ttfamily,
  breaklines       = true,
  captionpos       = b,
  frame            = single,
  rulecolor        = \color{codeframe},
  keepspaces       = true,
  numbers          = none,
  showspaces       = false,
}

\lstset{style=pramanix}

%% ── Hyperlinks ───────────────────────────────────────────────────────────────
\usepackage[hidelinks]{hyperref}

%% ── Headers and footers ──────────────────────────────────────────────────────
\usepackage{fancyhdr}
\pagestyle{fancy}
\fancyhf{}
\fancyhead[L]{\small\nouppercase{\leftmark}}
\fancyhead[R]{\small Pramanix v1.0.0}
\fancyfoot[C]{\thepage}
\renewcommand{\headrulewidth}{0.4pt}

%% ── Line spacing ─────────────────────────────────────────────────────────────
\usepackage{setspace}
\onehalfspacing

%% ── Paragraph style ──────────────────────────────────────────────────────────
\setlength{\parindent}{0pt}
\setlength{\parskip}{0.7em}

%% ── Lists ────────────────────────────────────────────────────────────────────
\usepackage{enumitem}

%% ── Custom shorthands ────────────────────────────────────────────────────────
\newcommand{\fname}[1]{\texttt{#1}}          % file / path name
\newcommand{\code}[1]{\texttt{#1}}           % inline code token
\newcommand{\cfrsec}[1]{\S\,#1}             % regulatory section

%% ── Title metadata ───────────────────────────────────────────────────────────
\title{%
  \Huge\bfseries PRAMANIX\\[0.6em]
  \Large A Deterministic Neuro-Symbolic Execution Firewall\\[0.3em]
  for Autonomous AI Agents\\[1.2em]
  \large A Technical Thesis:\\
  Engineering Design, Research Foundations, and Production Evidence
}
\author{Viraj Jain}
\date{Audit Date: 12 May 2026}

%% ============================================================================
\begin{document}
%% ============================================================================

\maketitle
\thispagestyle{empty}

\vfill
\begin{center}
  \small
  System Version: Pramanix 1.0.0 \quad\textbullet\quad License: AGPL-3.0-only\\[0.4em]
  Repository: \url{https://github.com/viraj1011JAIN/Pramanix}\\[1em]
  \textit{Evidence policy: every claim is grounded in a specific file, line number, or\\
  benchmark result from the actual codebase. Where a gap exists it is named exactly\\
  as a gap. Where the code proves a claim, the proof is cited.}
\end{center}

\newpage

%% ── Abstract ─────────────────────────────────────────────────────────────────
\chapter*{Abstract}
\addcontentsline{toc}{chapter}{Abstract}

Pramanix is a Python SDK that provides formal, deterministic verification of
AI agent actions using the Microsoft Research Z3 SMT solver.  Its core
contract is stated in \fname{docs/Blueprint.md} and enforced by code:
\code{Decision.allowed=True} if and only if Z3 returns \code{sat} for every
policy invariant against the given inputs.  Every other outcome ---
\code{unsat}, \code{unknown}, solver timeout, internal exception, Pydantic
validation failure, or network error --- produces \code{allowed=False}.
There is no third outcome and no code path in the codebase that returns
\code{allowed=True} through an error handler.

\textbf{The problem.}  Language models are stochastic.  When an AI agent
calls a bank transfer API, it reasons probabilistically over token sequences.
A system operating at 99.9\,\% accuracy across 10,000 daily decisions
misclassifies ten per day.  In financial and healthcare contexts, those ten
errors are not acceptable as a statistical artefact; they are individual
financial losses or clinical safety events.  Pramanix moves the accept/reject
decision from probabilistic inference to formal arithmetic verification.  It
does not make the LLM more accurate.  It enforces a hard envelope around what
the LLM's output can cause.

\textbf{What the codebase delivers} (verified against source as of
2026-05-12 in \fname{reality.md}):
\begin{itemize}
  \item A Python-native DSL that lowers to Z3 AST without \code{eval} or
        \code{exec} anywhere in the codebase (\fname{docs/DECISIONS.md}
        \cfrsec{9}).
  \item Exact rational arithmetic via \code{Decimal(str(v)).as\_integer\_ratio()}
        $\to$ \code{z3.RatVal()}, preventing IEEE~754 rounding errors in
        financial invariants (\fname{transpiler.py}, float handling path).
  \item Two-phase Z3 solving: a shared solver for the common SAT case, then
        per-invariant attribution solvers for UNSAT (\fname{docs/DECISIONS.md}
        \cfrsec{3}).
  \item Eleven-layer injection defence covering NFKC normalisation, regex
        pattern matching, length enforcement, and dual-LLM consensus
        (\fname{translator/}).
  \item Ed25519-signed decisions with Merkle batch anchoring and a
        self-contained offline verifier (\fname{audit/}).
  \item AES-256-GCM at-rest encryption for archive segments
        (\code{EncryptedArchiveWriter} in \fname{audit/archiver.py},
        Phase~10).
  \item Thirty-eight pre-built policy primitives across seven domains
        (\fname{reality.md} \cfrsec{2.2}).
  \item A compliance reporter mapping Z3 violation labels to regulatory
        citations in six legal frameworks, with PDF export
        (\fname{helpers/compliance.py}).  This is a tool for compliance
        engineers building structured evidence packages, not a legal opinion.
  \item Nine framework adapters (LangChain, LlamaIndex, AutoGen, FastAPI,
        CrewAI, DSPy, Haystack, Semantic Kernel, Pydantic AI).
        \textbf{Four of these --- DSPy, Haystack, Semantic Kernel, and
        Pydantic AI --- are stub implementations} as of v1.0.0, documented
        in \fname{docs/KNOWN\_GAPS.md}.
  \item A strict Intermediate Representation (IR) compiler
        (\fname{src/pramanix/compiler.py}) --- Pillar~1 of the
        Neuro-Symbolic Policy Engine --- that accepts \code{PolicyIR}, a
        Pydantic-typed JSON schema usable directly as an LLM
        \code{response\_format}, and lowers it deterministically to
        \code{ConstraintExpr} invariants identical to hand-authored ones.
        A companion \code{Decompiler} walks the compiled DSL AST and
        produces a timestamped structured-English audit report for CISO
        sign-off.  No LLM call is made by the compiler itself.
\end{itemize}

\textbf{Test suite state} (Python~3.13.7, pytest~8.4.2): \textbf{4,002 passed},
81 skipped, 0 failed.  Branch coverage: \textbf{98.26\,\%} (98\,\% gate:
passing).

\textbf{Self-audit findings} (\fname{reality.md}, 2026-05-12): The codebase
audit was conducted by the author against the author's own source code.  It
is not an independent third-party audit.  The structured rubric in
\fname{reality.md} produced a self-assigned score of 8.5/10.  The primary
adoption blocker is the AGPL-3.0-only license: any organization deploying
Pramanix as part of a network service must release the complete corresponding
source code of that service to users.  No commercial software company in
financial services or healthcare accepts this in a proprietary stack.  A
license change is the prerequisite for enterprise deployment.

This document is a technical record of what was built, what decisions were
made and why, where the gaps are, and what the evidence shows.  Where claims
are stronger than the evidence supports, that is stated explicitly.

\newpage
\tableofcontents
\newpage

%% ============================================================================
\chapter{Background and Motivation}
%% ============================================================================

\section{The Probabilistic Problem in AI Agents}

Modern AI agent frameworks --- LangChain, AutoGen, CrewAI, Semantic Kernel ---
give large language models tools: bank transfer APIs, infrastructure scale
commands, medical record access, trading order entry. The models are
stochastic. Their outputs are samples from a probability distribution. At
temperature zero they are more deterministic, but they are never formally
deterministic in the mathematical sense.

The practical consequence, articulated precisely in \fname{docs/Blueprint.md}
\cfrsec{1}, is this:

\begin{quote}
LLMs are stochastic token samplers. Temperature $> 0$ means no two identical
runs are guaranteed identical. In a banking context, a 99.9\,\% accuracy rate
on a financial operation means 1 in 1,000 transfers may be incorrect. At scale
--- 10,000 daily operations --- that is 10 incorrect decisions per day.
\end{quote}

This is not hypothetical. It is arithmetic: $10{,}000 \times 0.001 = 10$.

\section{Why Existing Approaches Fail}

The market's current answers fall into two broken categories, both documented
in the Blueprint.

\textbf{Category~1 --- Rule-Based Systems (regex, IF-THEN).}\ Cannot reason
about compound constraints. \code{amount < 10000 AND balance > amount AND NOT
frozen} looks reasonable until the edge cases arrive: floating-point rounding
on the balance comparison, compound conditions that interact non-linearly,
rules that were correct in isolation but produce wrong results in combination.
Rule-based systems have no completeness guarantee. A badly-written rule
silently passes or blocks incorrectly.

\textbf{Category~2 --- LLM-as-Judge.}\ Uses the same probabilistic tool to
judge itself. Adversarial prompts can override the judge entirely. The system
is circular.

Pramanix's approach, documented in \fname{docs/DECISIONS.md} \cfrsec{1},
rejects both:

\begin{quote}
For the decidable arithmetic theories Pramanix uses --- Linear Real
Arithmetic (LRA) and Linear Integer Arithmetic (LIA) --- Z3 either returns
\code{sat} (every invariant holds for the given inputs, with a satisfying
model as proof) or \code{unsat} (at least one invariant is violated, with an
unsatisfiable core for attribution).  Attribution via \code{unsat\_core()}
identifies exactly which invariants were violated.  Z3 handles real
arithmetic (\code{z3.Real}) natively, which is required for financial
invariants.  Z3 is deterministic: given the same formula and the same
\code{rlimit} cap, it produces the same answer.
\end{quote}

\textbf{Non-linear arithmetic caveat.}  For non-linear constraints (e.g.,
\code{E(price) * E(quantity) >= E(min\_total)}), Z3 may return \code{unknown}
rather than \code{sat} or \code{unsat}.  Pramanix is fail-closed:
\code{unknown} produces \code{allowed=False}.  Policy authors writing
multiplicative constraints should verify that Z3 returns \code{sat} for their
specific invariant under expected input ranges; a Z3 \code{unknown} result
produces a BLOCK decision, which is conservative but may be unexpected.

\section{The Name}

The name encodes the philosophy. \textit{Pram\={a}\d{n}a} is a Sanskrit
epistemological term meaning ``valid means of knowledge'' --- proof, evidence,
what makes a belief justified. \textit{Unix} is the composable systems
philosophy. Pramanix is the combination: a composable system that operates on
proof, not belief.

\section{What Pramanix Is Not}

The \fname{docs/Blueprint.md} \cfrsec{3} table prevents misuse:

\begin{table}[h]
\centering\small
\begin{tabular}{@{}p{4.4cm} p{9.2cm}@{}}
\toprule
\textbf{It is NOT} & \textbf{Because} \\
\midrule
A replacement for OPA &
  OPA handles AuthZ (who can try). Pramanix handles mathematical safety (is
  this attempt valid). Use both. \\[4pt]
A prompt guardrail &
  It does not filter LLM outputs. It verifies structured intents. \\[4pt]
A general-purpose rule engine &
  It operates exclusively on Z3-expressible constraints over typed fields. \\[4pt]
Secure without host freshness check &
  The \code{state\_version} check \emph{must} be implemented by the caller.
  Pramanix enforces the contract; it does not implement the check. \\
\bottomrule
\end{tabular}
\caption{What Pramanix is not (source: \fname{docs/Blueprint.md} \cfrsec{3})}
\end{table}

%% ============================================================================
\chapter{Research and Design Phase}
%% ============================================================================

\section{The Core Research Question}

The research question is practical and specific: can a Python library provide
formal, mathematical verification of AI agent actions with latency and
ergonomics suitable for production deployment?

Two sub-questions follow:
\begin{enumerate}
  \item Can Z3 be made fast enough for request-path verification?
  \item Can the policy authoring surface be ergonomic enough for application
        engineers who have never used an SMT solver?
\end{enumerate}

Both answers are yes, with caveats. The benchmark evidence is in
\fname{benchmarks/results/latency\_results.json}: P50 of 5.235\,ms, P95 of
6.361\,ms, P99 of 7.109\,ms over 2,000 samples in \code{sync} mode. The
ergonomics are answered by the DSL:
\code{(E(balance) - E(amount) >= 0).named("non\_negative\_balance")} is pure
Python, type-checked, IDE-complete, and readable without any knowledge of Z3.

\section{SMT Solvers as Verification Engines}

The choice of Z3 over alternative verification engines is documented in
\fname{docs/DECISIONS.md} \cfrsec{1} with reasons that are not post-hoc.

\textbf{OPA/Rego was considered and rejected.}\ OPA has built-in arithmetic
operators (\code{>}, \code{>=}, \code{+}, \code{-}, etc.) and can evaluate
numeric comparisons in Rego.  The reason for rejection is more specific:
OPA provides no formal satisfiability guarantee.  A Rego rule evaluating to
\code{false} means ``no derivation succeeded under the current rules'' ---
not ``this constraint is mathematically violated by the inputs.''  OPA has
no equivalent of Z3's \code{unsat\_core()}: identifying which of $k$ rules
caused a denial requires manual rule introspection or explicit trace
annotations, not a built-in attribution mechanism.  Additionally, OPA is
designed for authorization policy (who may perform what action), not
quantitative arithmetic safety verification (do the arithmetic relationships
between fields satisfy all constraints simultaneously?).  The recommended
deployment model is both: OPA for authorization, Pramanix for arithmetic
safety.

\textbf{Custom Python rule engine was considered and rejected.}\ Fast, but no
formal completeness guarantee. A badly-written rule silently passes or blocks
incorrectly with no attribution.

\textbf{Symbolic execution over Python bytecode was considered and rejected.}\
Complexity is prohibitive, failure modes are not well-understood.

Z3 wins on four properties: completeness guarantee (SAT means all invariants
hold), unsat-core attribution (exactly which invariants were violated), native
real arithmetic (exact, not floating-point), and determinism.

\section{The Exact Arithmetic Design Decision}

One design decision that separates Pramanix from naive Z3 usage deserves its
own section. \fname{docs/DECISIONS.md} \cfrsec{10} explains why floating-point
values are never passed directly to Z3:

\begin{quote}
Floating-point values in invariants are converted via
\code{Decimal(str(v)).as\_integer\_ratio()} before being passed to Z3 as
rational numbers (\code{z3.RatVal(numerator, denominator)}).
\end{quote}

The implementation is in \fname{transpiler.py} at the float handling path.
\code{Decimal(str(v)).as\_integer\_ratio()} produces an exact rational
(numerator, denominator pair). Z3's \code{RatVal} represents this as exact
rational arithmetic. $0.1 + 0.2 \neq 0.3$ in IEEE~754. In Z3 via this path,
it is exactly equal. Financial invariants of the form $\mathit{balance} -
\mathit{amount} \geq \mathit{min\_reserve}$ must be exact. This is not an
optimisation; it is a correctness requirement.

\section{The Two-Phase Solver Strategy}

The solver architecture in \fname{docs/DECISIONS.md} \cfrsec{3} reflects a
trade-off between the common case and the rare case.

\textbf{Phase~1 (fast path):}\ One shared \code{z3.Solver} with all
invariants added via \code{add()}. If \code{sat} $\Rightarrow$ return SAFE
immediately. This is the common case for legitimate requests.

\textbf{Phase~2 (attribution):}\ Each invariant gets its own \code{z3.Solver}
with \code{assert\_and\_track}. One assertion per solver so
\code{unsat\_core()} is always \code{\{label\}} --- complete violation
attribution with no minimal-core ambiguity.

The insight: \code{assert\_and\_track + unsat\_core()} on a single solver with
multiple assertions does not always return the minimal core --- Z3's core
extraction is heuristic. Using one solver per invariant guarantees that
\code{unsat\_core()} returns exactly \code{\{label\}}, so violation
attribution is complete and deterministic. The cost is $O(k)$ solves in
Phase~2, but Phase~2 only triggers on violations, which should be the minority
of requests in a correctly configured production system.

\section{The Fail-Closed Invariant}

The most important design decision in the system is documented in
\fname{docs/DECISIONS.md} \cfrsec{2}:

\begin{quote}
\code{Guard.verify()} catches all exceptions and returns
\code{Decision.error(allowed=False)}. It never raises to the caller.
\end{quote}

This is enforced not just by convention but by code structure.
\code{Decision.\_\_post\_init\_\_} raises \code{ValueError} if
\code{allowed=True} and \code{status != SAFE}. There is no code path in the
entire system where an error handler returns \code{allowed=True}. The
invariant is: Z3 returning SAT is the \emph{only} path to
\code{allowed=True}.

\textbf{Policy specification correctness.}  The fail-closed guarantee applies
to the policy as written.  Z3 proves that the given formula is satisfiable or
unsatisfiable for the given inputs.  It does not validate that the formula
correctly encodes the intended business rule.  A policy author writing
\code{E(amount) <= 10000} when they intended \code{E(amount) < 10000} will
receive a formally correct SAT result for \code{amount == 10000}, which is
business-wrong.  Policy specification review --- confirming that DSL
invariants accurately capture the intended business rules --- is a human
responsibility outside the scope of Z3.

%% ============================================================================
\chapter{Technical Architecture}
%% ============================================================================

\section{The Thirteen-Layer Architecture}

The system is not a single module. It is thirteen distinct layers with clean
boundaries, documented in \fname{docs/ARCHITECTURE\_NOTES.md} and
independently verified in \fname{reality.md} \cfrsec{1.2}.

\subsection{Layer~1 --- The Orchestrator (\fname{guard.py},
\fname{guard\_pipeline.py})}

The \code{Guard} class is the single entry point. Every \code{verify()} call
passes through a ten-step pipeline:

\begin{enumerate}[label=\arabic*.]
  \item Pydantic strict-mode validation (intent + state)
  \item Fast-path $O(1)$ pre-screen (if enabled) --- can only BLOCK
  \item Translator consensus extraction (if \code{translator\_enabled=True})
  \item \code{guard\_pipeline} semantic post-consensus check
  \item Z3 solve via \fname{solver.py}
  \item Decision construction
  \item Ed25519 signing (\code{\_sign\_decision})
  \item Audit sink emit (all configured sinks, failures caught and logged)
  \item Structlog emission
  \item OTel span close / Prometheus update
\end{enumerate}

The boundary invariant: \code{Guard.verify()} \textbf{never raises}. Every
exception, including unexpected ones from user-defined \code{invariants()}, is
caught and returned as \code{Decision.error()} with \code{allowed=False}.

\subsection{Layer~2 --- The DSL (\fname{expressions.py}, \fname{policy.py})}

Pure Python, zero Z3 imports. The \code{E()} builder, \code{ConstraintExpr},
\code{ForAll}, \code{Exists}, \code{ArrayField}, \code{DatetimeField},
\code{NestedField}. No string phase, no \code{eval}, no \code{exec} anywhere
in the codebase (\fname{docs/DECISIONS.md} \cfrsec{9}).

\code{ConstraintExpr.\_\_bool\_\_} raises \code{TypeError} rather than
returning a value (verified: \fname{expressions.py}:931). This is intentional:
\code{bool(expr1) and bool(expr2)} short-circuits at the Python level and
discards one operand. The policy author intended \code{expr1 \& expr2}.
Raising at compilation time surfaces the bug immediately rather than silently
producing wrong constraints. Both \code{ExpressionNode.\_\_bool\_\_} (line~787)
and \code{ConstraintExpr.\_\_bool\_\_} (line~931) raise \code{TypeError} with
an explanatory message directing authors to use \code{\&}/\code{|}/\code{\~{}}.

\subsection{Layer~2a --- Pillar~1 IR Compiler (\fname{src/pramanix/compiler.py})}

\fname{compiler.py} is the deterministic boundary between structured LLM
output and the Z3 verification engine.  It is distinct from
\fname{natural\_policy/} (Phase~2), which requires an active LLM client for
text-to-policy translation.  The IR compiler receives a Pydantic-validated
\code{PolicyIR} and lowers it to \code{ConstraintExpr} invariants that are
functionally identical to hand-authored ones --- the same objects
\fname{transpiler.py} receives from \code{Policy.invariants()}.

The module has four isolated layers:

\begin{enumerate}
  \item \textbf{IR schema.}  \code{FieldReference}, \code{Operator},
        \code{Condition}, \code{Rule}, \code{PolicyIR} --- pure Pydantic~v2,
        zero Z3, zero I/O.  \code{ConfigDict(extra="forbid", frozen=True)}
        on every model.  \code{Rule} is recursive (self-referential) via
        \code{Rule.model\_rebuild()}.  \code{PolicyIR} enforces globally
        unique rule names via a \code{model\_validator} --- duplicate names
        would produce duplicate invariant labels, which the Z3 attribution
        phase cannot disambiguate.

  \item \textbf{PolicyCompiler.}  Stateless;
        \code{compile(ir, policy\_cls) -> list[ConstraintExpr]}.  Validates
        at compile time: (a)~field existence against
        \code{Policy.fields()}; (b)~Z3 sort compatibility between \code{lhs}
        and a scalar \code{rhs}; (c)~field-to-field sort compatibility when
        \code{rhs} is also a \code{FieldReference}; (d)~ordering operators
        (\code{GT}, \code{LT}, \code{GTE}, \code{LTE}) rejected on
        \code{Bool}-sorted fields; (e)~\code{IN}/\code{NOT\_IN} requires a
        list RHS.  Every violation raises \code{PolicyCompilationError} ---
        partial compilation is architecturally impossible.  Uses
        \code{Decimal(str(scalar))} for \code{Real}-sorted scalars,
        consistent with \fname{transpiler.py}'s exact arithmetic path.

  \item \textbf{Fail-closed validation.}  Seven compile-time error classes
        are detected before any Z3 formula is constructed: unknown field;
        scalar type incompatible with field Z3 sort; field-to-field sort
        mismatch; ordering operator on \code{Bool} sort; membership operator
        on non-list RHS; empty membership set; \code{bool} literal as
        ordering RHS.

  \item \textbf{Decompiler.}  Reverse translation:
        $\code{list[ConstraintExpr]} \to \text{str}$.  Walks the internal
        DSL AST (\code{\_FieldRef}, \code{\_Literal}, \code{\_BinOp},
        \code{\_CmpOp}, \code{\_BoolOp}, \code{\_InOp}) without touching Z3.
        Produces a timestamped CISO sign-off report.  Unknown AST node types
        degrade gracefully to a \code{<TypeName>} placeholder --- the
        decompiler is intentionally forward-compatible.  Deterministic and
        idempotent.
\end{enumerate}

\textbf{Security invariants.}  No \code{eval()}, no \code{exec()}, no
dynamic code generation.  The LLM is never called by this module.  Every
field reference is validated against the \code{Policy} class's declared
\code{Field} attributes at compile time --- a hallucinated field name in the
LLM output is caught before any Z3 formula is constructed.  The \code{Operator}
and \code{FieldSource} enumerations are closed (\code{StrEnum} values) ---
an unknown operator string is rejected by Pydantic validation before the
compiler is invoked.

\textbf{Relationship to Phase~2 (\fname{natural\_policy/}).}  Phase~2's
\code{NaturalPolicyCompiler} is the management-plane component: it calls an
LLM and translates raw English text to a \code{NaturalPolicySchema}.
Pillar~1's \code{PolicyCompiler} is the data-plane compiler: it operates on
an already-structured \code{PolicyIR} with no LLM dependency.  The two can
be composed --- an intermediate adapter converts \code{NaturalPolicySchema}
to \code{PolicyIR} --- but each is independently usable.

\subsubsection{Polymorphic RHS and Pillar~1 + Pillar~2 Unification}

A \code{Condition} in the IR represents a single binary assertion
$\mathrm{lhs}\ \mathrm{op}\ \mathrm{rhs}$.  In earlier designs the
right-hand side was constrained to raw scalar literals, which prevented two
important classes of policy:

\begin{enumerate}
  \item \emph{Relational field-to-field assertions}: for example, asserting
    that a wire-transfer amount does not exceed the current account balance
    (\(\mathit{intent.amount} \leq \mathit{state.balance}\)) without
    hard-coding a numeric threshold.
  \item \emph{Identity-bound assertions}: checking the cryptographic identity
    of the calling agent (e.g.\
    $\mathit{intent.\_mesh\_principal} = \texttt{"spiffe://\dots/payments-agent"}$)
    where the identity value is a constant literal but the left-hand side
    is a field injected by \textbf{Pillar~2} (\code{MeshAuthenticator}).
\end{enumerate}

\textbf{Discriminated union.}  The \code{rhs} field of every \code{Condition}
is now a Pydantic~v2 discriminated union keyed by the JSON field \code{"type"}:

\begin{itemize}
  \item \code{FieldReference} (\code{type = "field"}) --- a field pointer
    into either the \code{intent} or \code{state} namespace, enabling
    field-to-field comparisons.
  \item \code{LiteralValue} (\code{type = "literal"}) --- a typed constant
    (\code{bool}, \code{int}, \code{float}, or \code{str}) or a non-empty
    list of scalars for \code{IN}/\code{NOT\_IN} membership tests.
\end{itemize}

Pydantic routes JSON to the correct branch before any application
validation runs, eliminating the ambiguous type-inference fallback that
previously caused silent coercion bugs.  A \code{model\_validator} enforces
structural constraints (e.g.\ \code{FieldReference} RHS is incompatible with
membership operators; boolean literals cannot appear with ordering operators).

\textbf{\code{\_mesh\_principal} bridge.}  \code{MeshAuthenticator.\allowbreak
authenticate\_and\_bind()} (Pillar~2) injects the SPIFFE URI of the verified
caller into the intent context under the key \code{"\_mesh\_principal"}.
Policy authors reference it as a \code{FieldReference} with
\code{source = "intent"} and \code{field\_name = "\_mesh\_principal"}.
Because \code{FieldReference.field\_name} accepts any non-empty string
(underscore prefixes are explicitly permitted, and are \emph{not} subject to
the snake\_case label-validation regex that applies to rule and condition
labels), this design requires no special-casing in the compiler or IR schema.
At compile time the field name is resolved against the target \code{Policy}
class's declared \code{Field} attributes; a hallucinated or missing field name
produces a \code{PolicyCompilationError} before any Z3 formula is constructed.

\textbf{Natural-language decompilation.}  The \code{\_OP\_PHRASE} module-level
dictionary maps each \code{Operator} to a human-readable verb phrase.  The
\code{Decompiler} uses this mapping for field-to-field conditions so that the
CISO audit report reads, for example:

\begin{quote}
\textit{``intent.amount must be less than or equal to state.balance''}
\end{quote}

\noindent rather than the raw operator symbol~\code{<=}.  This satisfies the
audit-readability requirement without requiring a separate NL generation model.

\textbf{Example.}  A complete policy YAML demonstrating all three RHS
variants --- \code{FieldReference}, scalar \code{LiteralValue}, and list
\code{LiteralValue} --- together with the \code{\_mesh\_principal} identity
gate is provided in \fname{examples/mesh\_policy\_ir.yaml}.

\subsection{Layer~3 --- The Transpiler (\fname{transpiler.py})}

The only file in the codebase that calls \code{z3.*} from DSL nodes. Lowers
Python AST to \code{z3.ExprRef} via explicit structural pattern matching
(\code{match}/\code{case}) starting at line~382.

Two critical details: floats go through \code{Decimal(str(v)).as\_integer\_ratio()}
before becoming \code{z3.RatVal}. Integer literals default to \code{z3.RealVal}
for compatibility with Real-sorted fields. \code{InvariantASTCache} pre-walks
the expression tree at Guard construction time --- no repeated tree walks at
request time.

\textit{Previously noted as a bug; verified fixed in code}: \code{\_tree\_repr()}
contains explicit \code{match}/\code{case} clauses for both \code{\_PowOp}
and \code{\_ModOp} (\fname{transpiler.py}:762--765). Polynomial and modulo
policies produce correct structural fingerprints. Z3 verification correctness
and policy diff tooling are unaffected.

\subsection{Layer~4 --- The Solver (\fname{solver.py})}

Z3 invocation and two-phase verification. Thread-local \code{z3.Context()} per
OS thread --- created once, never destroyed. This prevents GC races on
Windows/Python~3.13. Both wall-clock timeout
(\code{set("timeout", timeout\_ms)}) and elementary operation cap
(\code{set("rlimit", solver\_rlimit)}) are applied. The \code{rlimit} prevents
logic-bomb and non-linear expression DoS attacks that could exhaust CPU time
regardless of wall-clock timeout.

\subsection{Layer~5 --- The Worker Pool (\fname{worker.py})}

Three execution modes:

\begin{itemize}
  \item \code{sync} --- Z3 runs in the calling thread. Lowest latency, highest
        risk (Z3 crash kills the process).
  \item \code{async-thread} --- \code{ThreadPoolExecutor}. GIL is released
        during Z3 solving (Z3 releases the GIL), enabling genuine parallelism.
  \item \code{async-process} --- \code{ProcessPoolExecutor(mp\_context=spawn)}.
        Full process isolation. HMAC-sealed IPC prevents forged results from a
        compromised worker.
\end{itemize}

The \code{spawn} context (not \code{fork}) is mandatory because forking after
Z3 initialisation produces undefined behavior (\fname{docs/DECISIONS.md}
\cfrsec{4}). In \code{async-process} mode, each solve result is HMAC-SHA256
tagged with an ephemeral key (\code{\_EphemeralKey}) before crossing the
\code{multiprocessing.Queue} boundary. \code{\_EphemeralKey.\_\_reduce\_\_}
raises \code{TypeError} to prevent serialisation (\fname{docs/DECISIONS.md}
\cfrsec{15}).

\subsection{Layer~6 --- The Fast-Path Pre-Screener (\fname{fast\_path.py})}

\code{SemanticFastPath} and \code{FastPathEvaluator} --- configurable
pure-Python $O(1)$ rules that run after Pydantic validation and before Z3.
Architecture contract from \fname{docs/DECISIONS.md} \cfrsec{8}:

\begin{quote}
Fast-path rules can only BLOCK --- they return a string reason or \code{None}.
\code{None} passes through to Z3. There is no fast-path mechanism to allow a
request.
\end{quote}

Built-in rules: \code{negative\_amount}, \code{zero\_or\_negative\_balance},
\code{account\_frozen}, \code{exceeds\_hard\_cap},
\code{amount\_exceeds\_balance}.

\subsection{Layer~7 --- The LLM Translation Layer (\fname{translator/})}

Seven provider adapters: \code{openai\_compat}, \code{anthropic},
\code{ollama}, \code{gemini}, \code{cohere}, \code{mistral},
\code{llamacpp}. A five-layer injection defence stack executes before any LLM
call:

\begin{enumerate}
  \item Unicode NFKC normalisation --- collapses homoglyphs
  \item Input length enforcement --- \code{InputTooLongError} at
        \code{max\_input\_chars}
  \item Control-character stripping
  \item Injection pattern detection --- pre-compiled regex, $\sim$25 patterns
        covering instruction overrides, jailbreak keywords, open-source model
        tokens (Llama~2/3, ChatML, Phi-3), role escalation
  \item Dual-model consensus --- two independent LLMs must agree
\end{enumerate}

The dual-model consensus (\fname{docs/DECISIONS.md} \cfrsec{7}): a successful
prompt injection attack must manipulate both models in the same direction
simultaneously.  This is a harder attack than manipulating one model, provided
the two models are from different providers or different architectural families.
If both models share training data or come from the same provider family, a
shared adversarial example may fool both.  \textbf{Operational cost}: when
\code{translator\_enabled=True}, dual consensus means two LLM API calls per
verification request --- 2$\times$ API cost, 2$\times$ latency, and 2$\times$
rate-limit exposure.  This overhead is not captured in the benchmark tables,
which measure the Z3 verification path only.

\subsection{Layer~8 --- The Audit Chain (\fname{audit/})}

\code{DecisionSigner} signs \code{decision.decision\_hash} with Ed25519 via
\code{PramanixSigner}. \code{DecisionVerifier} is intentionally
self-contained --- an auditor can copy it as a single file and verify tokens
offline without installing Pramanix (\fname{docs/ARCHITECTURE\_NOTES.md}).

\code{MerkleAnchor} builds an in-memory Merkle tree:
$\code{add}(\mathit{decision\_id}) \to \code{root}() \to
\code{prove}(\mathit{decision\_id}) \to \code{proof.verify}()$. Proves any
single decision was part of an unaltered batch without replaying all decisions.

Archive encryption at rest (\textbf{resolved in Phase~10}):
\code{EncryptedArchiveWriter} (\fname{audit/archiver.py}) provides
AES-256-GCM encryption as a drop-in \code{archive\_writer=} pluggable
callback.  Each archive segment is encrypted with a fresh 96-bit nonce
(\code{secrets.token\_bytes(12)}) per NIST~SP~800-38D \cfrsec{8.2} and
written atomically via \code{tempfile.mkstemp()} + \code{os.replace()}.
GCM authentication tags verify byte-level integrity on every decrypt ---
any tampered byte raises \code{cryptography.exceptions.InvalidTag} before
any plaintext is exposed.  Key management stays with the caller (KMS/HSM).
The archiver now emits a \code{WARNING} \emph{only} when no
\code{archive\_writer=} is supplied.

The iterative Merkle root construction (replacing recursive) was a deliberate
correctness fix (\fname{docs/DECISIONS.md} \cfrsec{18}). The recursive
implementation hits Python's call stack limit at approximately 1,000 decisions.
The iterative version uses $O(1)$ stack depth.

\subsection{Layer~9 --- The Governance Gates}

\textbf{Information-Flow Control} (\fname{ifc/}): Lattice-based enforcement
between trust levels. \code{FlowPolicy} maps
\code{(source\_label, dest\_label)} pairs to \code{ALLOW} or \code{DENY}
rules. The \code{regulated()} preset explicitly rules that \code{REGULATED}
$\to$ \code{INTERNAL}, \code{REGULATED} $\to$ \code{CUSTOMER}, and
\code{REGULATED} $\to$ \code{CONFIDENTIAL} are all \code{permitted=False},
surfacing violation reasons in \code{FlowViolationError}.

\code{OversightWorkflow} (\fname{oversight/}) manages escalation queues.
\code{PrivilegeScope} (\fname{privilege/}) manages capability manifests and
execution scope boundaries. These gates run after consensus and before Z3.

\subsection{Layer~10 --- Compliance and Lifecycle Tools}

\code{ComplianceReporter} (\fname{helpers/compliance.py}) takes a
\code{Decision} object and produces structured regulatory citations: BSA/AML
(31~CFR~\S\,1020), OFAC/SDN (50~CFR~\S\,598), SEC wash sale (IRC~\S\,1091),
HIPAA (45~CFR~\S\,164), SOX (15~U.S.C.~\S\,7241), Basel~III (BCBS~189). JSON
and PDF export.  This reporter is a tool for compliance engineers building
structured evidence packages; it maps Z3 violation labels to cited regulations.
It is not a substitute for legal counsel or a compliance officer's judgment.

\code{PolicyDiff} and \code{ShadowEvaluator} (\fname{lifecycle/diff.py})
provide safe canary promotion. \code{ShadowEvaluator} runs the candidate
policy alongside every live decision, recording divergence events non-blocking
via \code{threading.Lock}.

\subsection{Layer~10a --- Pillar~3 Compliance Oracle
  (\fname{src/pramanix/compliance/oracle.py})}

\code{ComplianceOracle} is the regulatory attestation engine introduced in
Pillar~3.  It operates exclusively on completed \code{ProvenanceRecord}
objects and is \emph{structurally decoupled} from \code{Guard.verify()}: it
imports nothing from the hot-path layers (\fname{guard.py},
\fname{solver.py}, \fname{transpiler.py}).  Its sole function is to read a
finished audit record and emit a cryptographically-bound
\code{ComplianceAttestation} mapping the record to one or more regulatory
control identifiers.

\subsubsection{Design Invariant: Offline-Only Evaluation}

The oracle is intentionally an \emph{offline} component.  Coupling compliance
attestation to the synchronous verification hot-path would (a)~add variable
latency to safety-critical decisions, (b)~create a single point of failure
that could, if the oracle raised, cause the guard to block on an internal
error rather than on a policy violation.  Instead, the oracle is called
asynchronously after the guard has already written the \code{ProvenanceRecord}
to the audit sink.  This matches the architectural pattern used by hardware
security modules for post-hoc attestation.

\subsubsection{\texttt{ControlMapping} and Discriminated Evidence}

A \code{ControlMapping} is a frozen Pydantic~v2 model with two optional
evidence criteria:

\begin{itemize}
  \item \code{invariant\_label} (\code{str | None}) --- matched by exact
        equality against invariant labels in \code{evaluated\_invariants}
        or \code{violated\_invariants}.
  \item \code{principal\_pattern} (\code{str | None}) --- matched via
        \code{fnmatch.fnmatch} against \code{record.principal\_id}, which
        carries the verified SPIFFE URI injected by Pillar~2's
        \code{MeshAuthenticator}.
\end{itemize}

At least one criterion must be set (enforced in \code{model\_post\_init});
\code{require\_both=True} (default) requires both criteria to match when
both are specified.  \code{require\_both=False} produces OR semantics.  The
\code{MappingMatchKind} enum (\code{INVARIANT\_LABEL},
\code{PRINCIPAL\_IDENTITY}, \code{BOTH}) records which criterion triggered.

This discriminated evidence model directly bridges Pillar~1 (invariant
labels produced by the IR compiler and Z3 attribution phase) with Pillar~2
(SPIFFE identity injected by \code{MeshAuthenticator}).  A single
\code{ControlMapping} can therefore express a compound regulatory assertion
such as: ``this principal identity must have invoked a tool that satisfied
this invariant'' --- with no additional glue code.

\subsubsection{Cryptographic Provenance in \texttt{ComplianceAttestation}}

Every \code{ComplianceAttestation} embeds a \code{record\_hmac\_tag} field
derived from \code{ProvenanceRecord.hmac\_tag()}.  This HMAC-SHA256 tag is
computed over the immutable fields of the originating record, providing a
cryptographic link from the attestation back to the exact record that was
evaluated.  The oracle also accepts a \code{stored\_hmac\_tag} parameter
for cross-process audit scenarios in which the live record object is
unavailable --- the stored tag from the original write can be supplied
directly.

\subsubsection{Fail-Closed Evaluation}

\code{evaluate\_record()} has a single, unbreachable contract:
it \emph{never raises}.  The outermost try-except catches all exceptions,
logs them at \code{ERROR} level, and returns a minimal
\code{ComplianceAttestation} with \code{total\_controls\_matched = 0} and a
summary string recording the failure.  No unhandled exception in the oracle
can propagate to the caller and disrupt audit pipeline execution.

\subsubsection{Thread Safety}

The internal registry (\code{defaultdict(list)}) is protected by a
\code{threading.RLock}.  \code{register\_mapping()} holds the lock for the
duration of the append.  \code{evaluate\_record()} snapshots the registry
under lock into a local \code{dict[RegulatoryFramework, list[ControlMapping]]}
and immediately releases the lock before any evaluation work begins.  This
design allows the oracle to serve concurrent audit pipeline workers without
holding the lock across (potentially slow) evaluation loops.

\subsubsection{Public Surface}

Eight symbols are re-exported from \code{pramanix.compliance} and from the
top-level \code{pramanix} package: \code{ComplianceAttestation},
\code{ComplianceOracle}, \code{ControlEnforcementResult},
\code{ControlMapping}, \code{ControlSatisfactionResult},
\code{FrameworkAttestation}, \code{MappingMatchKind},
\code{RegulatoryFramework}.

\subsection{Layer~11 --- The Request Infrastructure}

\code{ResolverRegistry} stores per-request resolved field values in a
\code{contextvars.ContextVar} --- not \code{threading.local}. The distinction
is security-critical: in asyncio, multiple concurrent tasks share one OS
thread. \code{threading.local} would allow Task~B to see Task~A's resolved
values --- a P0 data bleed between users (\fname{docs/DECISIONS.md}
\cfrsec{6}). \code{Guard} calls \code{clear\_cache()} in its \code{finally}
block unconditionally --- no resolved value survives across requests.

\code{validator.py} uses Pydantic~v2 strict mode. Implicit coercions
(\code{"123"} $\to$ \code{int}) are rejected at this single boundary.

\subsection{Layer~12 --- The Exception Hierarchy (\fname{exceptions.py})}

Nineteen typed exception classes in a strict \code{PramanixError} tree.
Domain-specific: \code{GuardViolationError} (contains the full \code{Decision}
object), \code{FlowViolationError}, \code{PrivilegeEscalationError},
\code{OversightRequiredError}, \code{MemoryViolationError},
\code{ProvenanceError}, \code{IntegrityError}, \code{SolverTimeoutError},
\code{ExtractionMismatchError}. Callers never need to catch bare
\code{Exception}.

\subsection{Layer~13 --- The Primitives Library (\fname{primitives/})}

Thirty-eight pre-built Policy constraints across seven domains
(\fname{reality.md} \cfrsec{2.2}):

\begin{itemize}
  \item \textbf{Finance:} \code{NonNegativeBalance}, \code{UnderDailyLimit},
        \code{UnderSingleTxLimit}, \code{RiskScoreBelow}
  \item \textbf{FinTech:} \code{AntiStructuring}, \code{CollateralHaircut},
        \code{KYCTierCheck}, \code{MarginRequirement}, \code{MaxDrawdown},
        \code{SanctionsScreen}, \code{SufficientBalance},
        \code{TradingWindowCheck}, \code{VelocityCheck}, \code{WashSaleDetection}
  \item \textbf{Healthcare:} \code{BreakGlassAuth}, \code{ConsentActive},
        \code{DosageGradientCheck}, \code{PediatricDoseBound},
        \code{PHILeastPrivilege}
  \item \textbf{RBAC:} \code{ConsentRequired}, \code{DepartmentMustBeIn},
        \code{RoleMustBeIn}
  \item \textbf{Infrastructure:} \code{BlastRadiusCheck},
        \code{CircuitBreakerState}, \code{CPUMemoryGuard}, \code{MaxReplicas},
        \code{MinReplicas}, \code{ProdDeployApproval}, \code{ReplicaBudget},
        \code{WithinCPUBudget}, \code{WithinMemoryBudget}
  \item \textbf{Time:} \code{After}, \code{Before}, \code{NotExpired},
        \code{WithinTimeWindow}
  \item \textbf{Common:} \code{FieldMustEqual}, \code{NotSuspended},
        \code{StatusMustBe}, \code{EnterpriseRole}, \code{HIPAARole}
\end{itemize}

\section{The Decision Object}

\code{Decision} is a frozen dataclass with one invariant enforced in
\code{\_\_post\_init\_\_}: $\code{allowed=True} \Leftrightarrow
\code{status=SAFE}$. No other combination is valid. The status values and
their semantics are shown in Table~\ref{tab:decision-status}.

\begin{table}[h]
\centering\small
\begin{tabular}{@{}l l p{7.5cm}@{}}
\toprule
\textbf{Status} & \textbf{allowed} & \textbf{Cause} \\
\midrule
\code{SAFE}                & \code{True}  & Z3 proved all invariants hold \\
\code{UNSAFE}              & \code{False} & Z3 found a counterexample \\
\code{TIMEOUT}             & \code{False} & Z3 exceeded time budget \\
\code{ERROR}               & \code{False} & Unexpected internal error \\
\code{STALE\_STATE}        & \code{False} & \code{state\_version} field mismatch \\
\code{VALIDATION\_FAILURE} & \code{False} & Pydantic model validation failed \\
\code{CONSENSUS\_FAILURE}  & \code{False} & Dual-model LLM disagreement \\
\code{RATE\_LIMITED}       & \code{False} & Load-shedding threshold exceeded \\
\bottomrule
\end{tabular}
\caption{Decision status values and their semantics}
\label{tab:decision-status}
\end{table}

The canonical hash is SHA-256 over \code{decision\_id + status + allowed +
violated\_invariants + explanation}, computed via \code{orjson} (or stdlib
\code{json} fallback). The \code{iat} timestamp is excluded from the signed
payload --- signatures are deterministic regardless of signing time
(\fname{docs/DECISIONS.md} \cfrsec{16}).

\section{The Public API Stability Contract}

Stability tiers from \fname{docs/PUBLIC\_API.md}:

\begin{itemize}
  \item \textbf{Stable:} \code{core}, \code{audit}, \code{crypto},
        \code{circuit\_breaker}, \code{execution\_token}, \code{key\_provider},
        \code{compliance}, \code{audit\_sinks}, \code{worker}, \code{primitives}
  \item \textbf{Beta:} \code{translator}, \code{integrations},
        \code{fast\_path}, \code{ifc}, \code{privilege}, \code{oversight},
        \code{memory}, \code{lifecycle}, \code{provenance}
\end{itemize}

Beta subsystems are usable in production but may change in minor versions with
deprecation notice.

%% ============================================================================
\chapter{Development Journey}
%% ============================================================================

\section{From Concept to 1.0.0}

The development history is reconstructed from \fname{docs/CHANGELOG.md} (Keep
a Changelog format) and the \fname{docs/Blueprint.md} milestone sequence.

\subsection{Phase~0 --- v0.0: Concept Validation}

The fundamental question: can Z3 verify a financial invariant in Python within
acceptable latency? Before any architecture was built, this had to be answered
empirically. The answer is yes --- the benchmark in
\fname{benchmarks/latency\_benchmark.py} demonstrates P99 of 7.109\,ms over
2,000 samples.

\subsection{Phase~1 --- v0.1--v0.3: Core Engine}

Foundation: \fname{guard.py}, \fname{policy.py}, \fname{expressions.py},
\fname{transpiler.py}, \fname{solver.py}. The key design decision made at
this stage was the fail-closed contract. \code{Decision.\_\_post\_init\_\_}
raises \code{ValueError} if \code{allowed=True} and \code{status != SAFE}.
This constraint is checked in every test suite run and has never been relaxed.

\subsection{Phase~2 --- v0.4: Exception Hierarchy and Type Safety}

Nineteen typed exceptions replacing bare raises. \code{mypy} strict mode
(\fname{pyproject.toml}: \code{mypy.strict = true}). The policy of zero
\code{cast()} except where Z3's incomplete type stubs require it
(\fname{transpiler.py}:390--396) was established here.

\subsection{Phase~3 --- v0.5: Security Hardening}

The formal threat model was written. HMAC-sealed IPC for \code{async-process}
mode (\fname{docs/DECISIONS.md} \cfrsec{15}) was implemented here.
\code{\_EphemeralKey.\_\_reduce\_\_} raising \code{TypeError} to prevent
serialisation was a direct response to the threat of worker process
exploitation.

The musl/Alpine rejection at import time (\fname{\_platform.py}) was added.
From \fname{docs/DECISIONS.md} \cfrsec{12}: Z3's glibc-compiled wheels
segfault or run 3--10$\times$ slower on musl. By the time Guard is
constructed, Z3 is already loaded. Failing at import time is the only way to
produce an attributable error in container startup logs.

\subsection{Phase~4 --- v0.6: Domain Primitives}

Thirty-eight pre-built primitives across seven domains. The
\code{ComplianceReporter} with six regulatory frameworks and PDF export. This
phase transformed the SDK from a framework requiring users to write all
constraints from scratch into a library with batteries included.

\subsection{Phase~5 --- v0.7: Performance Engineering}

The \code{InvariantASTCache} eliminates repeated tree walks at request time.
The \code{SemanticFastPath} provides $O(1)$ Python pre-screening before Z3.
The process pool with zero-IPC worker architecture, documented in
\fname{benchmarks/100m\_worker\_fast.py}:

\begin{lstlisting}[style=plain,caption={IPC architecture comparison (source:
  \fname{benchmarks/100m\_worker\_fast.py})}]
Old: asyncio coroutine -> ThreadPoolExecutor -> Guard process -> Z3 process
     2 IPC crossings per decision x ~15 ms each = ~30 ms overhead
     = 22-27 RPS/worker

New: 18 OS processes, each owns sync Guard + Z3 in one address space
     0 IPC crossings per decision
     = 80-120 RPS/worker  ->  1,440-2,160 aggregate RPS
\end{lstlisting}

Three micro-optimisations:
\begin{enumerate}
  \item \code{max\_input\_bytes=0} --- skips the per-decision JSON size-check
        round-trip ($\sim$1--2\,ms saved per decision)
  \item Pre-built payload cache --- 10,000 payloads generated at startup,
        no \code{Decimal()} construction in the decision loop ($\sim$0.5\,ms
        saved)
  \item \code{orjson} $+$ 8\,MiB buffer --- one OS \code{write()} per
        $\sim$100k decisions vs.\ one per decision ($\sim$0.3\,ms saved)
\end{enumerate}

\subsection{Phase~6 --- v0.8: Cryptographic Audit Engine}

Ed25519 signing (\fname{crypto.py}), \code{MerkleAnchor} and
\code{PersistentMerkleAnchor} (\fname{audit/merkle.py}),
\code{ProvenanceChain} (\fname{provenance.py}). The iterative Merkle root was
a correctness fix --- the recursive implementation hit Python's stack limit at
1,000 decisions (\fname{docs/DECISIONS.md} \cfrsec{18}).

\subsection{Phase~7 --- v0.9: Enterprise Integrations}

Nine framework adapters: LangChain, LlamaIndex, AutoGen, FastAPI, CrewAI,
DSPy, Haystack, Semantic Kernel, Pydantic AI. Two runtime interceptors: gRPC
\code{ServerInterceptor} and Kafka consumer wrapper. Kubernetes
\code{ValidatingWebhook}. The \code{@guard} decorator.

\textbf{Security bug (H-03)}: \code{PramanixGuardedTool} previously defaulted
\code{execute\_fn} to \code{lambda i: "OK"} when no function was provided.
Every guarded action returned success without executing any real logic. Fixed:
\code{execute\_fn=None} now raises \code{NotImplementedError}
(\fname{CHANGELOG.md} \cfrsec{H-03}).

\textbf{Security bug (H-04)}: \code{PramanixCrewAITool.\_run()} previously
returned a non-exception string when \code{underlying\_fn=None}. A blocked
CrewAI tool should raise, not return (\fname{CHANGELOG.md} \cfrsec{H-04}).

\subsection{Phase~8 --- v0.9.5: Security Hardening Pass}

\begin{itemize}
  \item \textbf{C-01}: \code{verify\_async()} had \code{except Exception: pass}
        in the \code{max\_input\_bytes} size-check path. Unserializable payloads
        silently bypassed the size gate. Fixed: returns
        \code{Decision.error(allowed=False)}.

  \item \textbf{H-02/M-05}: Timing pad applied only to BLOCK responses. ALLOW
        responses returned immediately, leaking a timing oracle. Fixed: timing
        pad applied unconditionally before the ALLOW/BLOCK branch.

  \item \textbf{BUG-03}: \code{OversightRecord} had no \code{\_\_slots\_\_}.
        Arbitrary attribute injection on a tamper-evident audit record. Fixed.

  \item \textbf{BUG-10}: \code{JWTIdentityLinker.\_verify\_token()} decoded the
        JWT payload before validating the \code{alg} header. An attacker could
        craft a token with \code{alg: "RS256"} and an HMAC-SHA256 signature,
        bypassing the asymmetric-key check (CVE-2015-9235 family). Fixed:
        \code{alg} validated as \code{"HS256"} \emph{before} signature
        computation.

  \item \textbf{BUG-11}: JWT \code{nbf} (not-before) claim not enforced.
        Tokens with future \code{nbf} were accepted. Fixed.

  \item \textbf{BUG-12}: Empty or missing JWT \code{sub} claim accepted.
        \code{"sub": ""} produced an empty-identity principal. Fixed.
\end{itemize}

\subsection{Phase~9 --- v1.0.0~GA: Documentation, Coverage, and Honest Gaps}

\fname{docs/KNOWN\_GAPS.md} was created --- an honest inventory of 14 known
limitations. \fname{MIGRATION.md} covering v0.7--v1.0 upgrade paths.
\fname{docs/PUBLIC\_API.md} stability tiers finalised. \code{pramanix doctor}
check~\#10 (logging-handlers) and check~\#11 (policy-hash-binding) added.

Test suite at v1.0.0 GA + Phase~10 hardening: 3,916 passed, 81 skipped, 0
failed.  Coverage: 98.16\,\% (98\,\% gate: \textbf{passing}).  Python~3.13.7,
pytest~8.4.2.  Total wall time: 417\,s on a single-core Windows CI node.

After Phase~11 test-isolation fixes: 3,916 passed, 81 skipped, 0 failed.
Coverage: 98.22\,\% (gate: \textbf{passing}).  Wall time: 793\,s.

After Phase~12 metric-registry and coverage fixes: 3,923 passed, 81 skipped,
0 failed.  Coverage: \textbf{98.26\,\%} (gate: \textbf{passing}).  Wall
time: 867\,s.

After Phase~2 (Natural Policy Engine) additions and test-suite bug fixes:
\textbf{3,925 passed}, \textbf{120 skipped}, \textbf{0 failed}.  Coverage:
\textbf{98.26\,\%} (gate: \textbf{passing}).  Wall time: $\approx$\,430\,s.

After pytest upgrade (7.4.4~$\to$~8.4.2) and full re-run (2026-05-13):
\textbf{4,002 passed}, \textbf{81 skipped}, \textbf{0 failed}.  Coverage:
\textbf{98.26\,\%} (gate: \textbf{passing}).  Wall time: 437.91\,s.

% ─────────────────────────────────────────────────────────────────────────────
\subsection*{Phase~2 (Natural Policy Engine) --- v1.1: Natural Language Policy Compiler}
% ─────────────────────────────────────────────────────────────────────────────

\paragraph{Motivation.}
Every prior policy in Pramanix was authored as Python code (decorator calls or
\code{PolicyConfig} objects).  That gates governance on engineering bandwidth:
a compliance officer cannot write a policy without knowing Z3's S-expression
vocabulary or Pramanix's \code{ConstraintExpr} DSL.  Phase~2 removes that
dependency by introducing a \emph{compile-time} natural-language front-end
that converts plain English sentences into verified, deterministic Z3
constraint graphs — with zero LLM inference on the hot evaluation path.

\paragraph{Five Non-Negotiable Design Constraints.}

\begin{enumerate}
\item \textbf{No LLM on the hot path.}  Every policy evaluation is purely
  deterministic.  An LLM is permitted only at \emph{compile time} (the
  author's workstation), never inside \code{Guard.verify()}.
\item \textbf{Full type safety.}  All intermediate representations are
  validated Pydantic v2 models; no raw dictionaries cross API boundaries.
\item \textbf{Meta-verification.}  Every compiled policy is checked with Z3
  for satisfiability and monotonicity before it is returned to the caller.
\item \textbf{Fail-closed.}  A malformed or unverifiable natural-language
  policy raises a hard error; it never silently permits a request.
\item \textbf{Python~3.11+.}  All type annotations use modern union syntax
  (\code{X | Y}) and pattern-matching where appropriate.
\end{enumerate}

\paragraph{Compilation Pipeline.}

\[
  \text{English sentence}
  \;\xrightarrow{\text{LLM (compile-time only)}}\;
  \code{NaturalPolicySchema}
  \;\xrightarrow{\text{ASTBuilder}}\;
  \code{ConstraintExpr}~(\text{Z3 DSL})
  \;\xrightarrow{\text{MetaVerifier}}\;
  \code{CompiledPolicy}
\]

\noindent
The four stages are logically independent modules; each can be replaced or
extended without touching the others.

\begin{description}

\item[\code{NaturalPolicySchema}]
  A Pydantic v2 model produced by the LLM from the natural-language input.
  Top-level fields: \code{policy\_name} (short human-readable label),
  \code{original\_english} (verbatim CISO text used by the meta-verifier),
  \code{fields} (list of \code{FieldDeclaration} each carrying \code{name},
  \code{z3\_type} $\in$ \{\code{Real}/\code{Int}/\code{Bool}/\code{String}\},
  and \code{description}), and \code{constraints} (a discriminated-union list
  of \code{comparison}, \code{and}, \code{or}, \code{not} nodes keyed by
  \code{kind}).  All fields are Pydantic-validated at construction time;
  constraint labels must match \code{\^{}[a-z][a-z0-9\_]*\$} (snake\_case)
  to serve as valid Python identifiers and prevent label injection.

\item[\code{ASTBuilder}]
  A stateless transformer that converts a \code{NaturalPolicySchema} into a
  list of \code{ConstraintExpr} objects using the existing Z3 DSL.  The
  module-level \code{\_Z3\_TO\_PYTHON} map converts each
  \code{FieldDeclaration.z3\_type} (a \code{Z3TypeEnum} value) to the
  corresponding Python type: \code{Real}→\code{Decimal}, \code{Int}→\code{int},
  \code{Bool}→\code{bool}, \code{String}→\code{str}.  The builder enforces
  type-operator legality: arithmetic operators (\code{+/-/*//}) are rejected
  on \code{Bool}/\code{String} fields, and ordering operators
  (\code{<}, \code{<=}, \code{>}, \code{>=}) are similarly rejected.  Any
  violation raises \code{FieldTypeError} before any Z3 context is entered.

\item[\code{NaturalPolicyCompiler}]
  The public entry point.  Given a \code{NaturalPolicySchema} (or a
  dictionary that Pydantic can coerce to one), it calls \code{ASTBuilder}
  to produce a \code{ConstraintExpr}, passes the result through
  \code{MetaVerifier}, and returns an immutable \code{CompiledPolicy} if
  verification succeeds.  Raises \code{ValueError} (schema) or
  \code{RuntimeError} (Z3 refutation) on any failure.

\item[\code{MetaVerifier}]
  Performs \emph{structural} (not Z3-semantic) anti-hallucination checks on
  the compiled AST.  For each constraint node the verifier reconstructs a
  canonical algebraic expression (e.g.\ \code{amount <= 50000}) from the
  compiled \code{ConstraintExpr} and compares it token-by-token against the
  LLM's \code{natural\_language} annotation, checking field names, operator
  synonyms (``at least'', ``not exceed'', etc.), and numeric values.  A
  mismatch indicates either an LLM hallucination (wrong operator or threshold)
  or a compiler bug (operator transposition).  Behaviour is governed by
  \code{VerificationMode}: \code{STRICT} raises \code{PolicyCompilationError},
  \code{WARN} records mismatches in \code{MetaVerificationResult.mismatches}
  without raising, and \code{SKIP} bypasses all checks.

\item[\code{CompiledPolicy}]
  A frozen dataclass with four attributes: \code{fields}
  (\code{dict[str, Field]} mapping each declared field name to its
  \code{Field} descriptor), \code{constraints}
  (\code{list[ConstraintExpr]} ready to be returned from
  \code{Policy.invariants()}), \code{verification}
  (\code{MetaVerificationResult} — full audit trail of the meta-verification
  pass), and \code{schema} (the \code{NaturalPolicySchema} intermediate
  representation produced by the LLM).  Attempting to mutate any field raises
  \code{FrozenInstanceError}.

\end{description}

\paragraph{Implementation Details.}

\begin{itemize}

\item \textbf{Module layout.}
  \fname{src/pramanix/natural\_policy/\_\_init\_\_.py} exports the full public
  surface: \code{NaturalPolicyCompiler}, \code{ASTBuilder},
  \code{CompiledPolicy}, \code{MetaVerifier}, \code{MetaVerificationResult},
  \code{VerificationMode}.  Internal helpers live in \fname{compiler.py},
  \fname{verifier.py}, and \fname{schemas.py} — none are exposed at the
  package root.

\item \textbf{Z3 thread safety.}  All Z3 solver operations inside
  \code{MetaVerifier} use a thread-local Z3 context following the established
  Phase~5 pattern, ensuring that concurrent Guard instances do not share
  solver state.

\item \textbf{No runtime LLM dependency.}  The compiler, builder, and
  verifier have zero imports from any LLM library.  The only optional
  dependency for end-to-end use is whatever translator is chosen to produce
  the \code{NaturalPolicySchema} from raw text — and that call happens
  entirely outside the compiler.

\end{itemize}

\paragraph{Test Coverage.}

Phase~2 ships \fname{tests/unit/test\_natural\_policy.py} — 61 new unit tests
across 9 test classes:

\begin{enumerate}
\item \code{TestZ3ToPythonMap} (5 tests) — verifies that all four
  \code{Z3TypeEnum} entries (\code{Real}, \code{Int}, \code{Bool},
  \code{String}) are present and map to the correct Python types (\code{Decimal},
  \code{int}, \code{bool}, \code{str}).
\item \code{TestNaturalPolicySchema} (8 tests) — Pydantic v2 validation:
  valid schema construction, duplicate field-name rejection, undeclared field
  reference in a constraint, invalid snake\_case label, \code{Bool} RHS with
  an arithmetic operator, dotted field names, empty \code{policy\_name}.
\item \code{TestASTBuilder} (14 tests) — comparison operators
  (\code{<=}/\code{>=}/\code{>}/\code{<}), \code{AND}/\code{OR}/\code{NOT}
  compositions, \code{ArithmeticLHS} (\code{balance - amount >= 0}),
  \code{FieldTypeError} on arithmetic with Bool/String fields,
  unknown-field rejection, field registry property, and multi-constraint
  parallel annotation.
\item \code{TestMetaVerifier} (10 tests) — \code{STRICT} passes on clean
  constraint; \code{SKIP} bypasses without calling helpers; \code{WARN}
  records mismatches in \code{MetaVerificationResult.mismatches} without
  raising; \code{STRICT} raises on hallucinated value; mismatched constraint
  lengths; per-constraint \code{ConstraintVerificationDetail} population;
  private helpers \code{\_get\_op\_synonyms} (``at least'',
  ``not exceed'') and \code{\_tokenise} (stop-word removal).
\item \code{TestNaturalPolicyCompiler} (10 tests, \code{async}) — async
  compilation with a mocked \code{Translator}: happy path,
  \code{ExtractionFailureError} on invalid JSON,
  \code{PolicyCompilationError} on undeclared field, \code{FieldTypeError}
  on type mismatch, \code{STRICT} verification integration, multiple
  field types and boolean combinators.
\item \code{TestCompileFromSchema} (3 tests) — synchronous
  \code{NaturalPolicyCompiler.compile\_from\_schema()} path with and without
  verification; \code{FieldTypeError} propagation.
\item \code{TestCompiledPolicy} (4 tests) — frozen dataclass mutation raises
  \code{FrozenInstanceError}; \code{verification} is a
  \code{MetaVerificationResult}; \code{schema} is a
  \code{NaturalPolicySchema}; \code{fields} contains \code{Field} instances.
\item \code{TestVerificationMode} (2 tests) — all three enum values
  (\code{STRICT}/\code{WARN}/\code{SKIP}) exist and are \code{str}
  subclasses.
\item \code{TestVerifierHelpers} (5 tests) — \code{\_reconstruct\_constraint}
  round-trips \code{>=}/\code{<=} and \code{ArithmeticLHS} nodes;
  \code{\_tokenise} removes stop words and is case-insensitive.
\end{enumerate}

All 61 tests pass in isolation and as part of the full suite.

\paragraph{Compile-Policy CLI Subcommand.}
Phase~2 adds a \code{compile-policy} subcommand to the \code{pramanix} CLI
(\fname{src/pramanix/cli.py}):

\begin{verbatim}
pramanix compile-policy POLICY_FILE [--verify {strict,warn,skip}] [--json]
pramanix compile-policy --example
\end{verbatim}

The subcommand reads a YAML or JSON file conforming to
\code{NaturalPolicySchema}, validates it with
\code{NaturalPolicySchema.model\_validate()}, compiles it via
\code{NaturalPolicyCompiler.compile\_from\_schema()}, and prints a
human-readable (or JSON) summary.  The \code{--example} flag prints a
self-contained YAML template to stdout and exits~0, allowing users to
bootstrap a schema without consulting the documentation.

\begin{itemize}
  \item \textbf{Verification mode}: \code{--verify strict} exits~0 if and
        only if \code{MetaVerificationResult.verified} is \code{True};
        \code{--verify warn} (default) exits~0 on any successfully compiled
        policy, printing mismatch warnings to stderr;
        \code{--verify skip} bypasses all meta-verification.
  \item \textbf{JSON output}: \code{--json} emits a structured object with
        keys \code{ok}, \code{policy\_name}, \code{rules\_compiled},
        \code{verification\_mode}, \code{verification\_passed}, and
        \code{verification\_warnings}.  Compilation errors also surface as
        JSON (\code{ok: false}) for machine-readable pipeline consumption.
  \item \textbf{Exit codes}: 0~(compiled + verified), 1~(compilation or
        verification failure), 2~(usage error — no file, no \code{--example}).
\end{itemize}

Because \code{compile\_from\_schema()} requires no LLM call, the subcommand
is suitable for CI pipelines that pre-validate natural-language policy files.
Eight unit tests in \fname{tests/unit/test\_cli\_compile\_policy.py} cover:
\code{--example} output, valid YAML compilation, \code{--json} output format,
all three verification modes, missing-file error, no-argument error, and JSON
error output on a malformed schema.

\paragraph{Injection Filter Hardening.}
Six new pattern tuples were added to
\fname{src/pramanix/translator/\_injection\_patterns.py}, covering two attack
classes absent from the original pattern set:

\begin{itemize}
  \item \textbf{\code{constraint\_override}} (5 patterns) — phrases that
        instruct the LLM to discard, ignore, override, or forget its operative
        constraints: \emph{discard (all) (previous) constraints/rules},
        \emph{ignore (all) constraints/policies/limits},
        \emph{override (all) constraints/policies/rules},
        \emph{forget (the) constraints/rules/guidelines},
        \emph{no (more) restrictions/constraints/limits/rules/policies}.
  \item \textbf{\code{policy\_bypass}} (1 pattern) — phrases explicitly
        requesting bypass of the governing policy or safety filters:
        \emph{bypass (the) policy/policies/filter/filters/constraints/safety}.
\end{itemize}

All six patterns use IGNORECASE-anchored regular expressions compiled once at
import time.  Ten new tests were added to \code{TestInjectionFilter} in
\fname{tests/unit/test\_injection\_scorer\_filter.py}, bringing the class total
to 28 tests, all passing.

A security comment was also added to the \code{pickle.loads(raw)} call in
\fname{src/pramanix/translator/injection\_scorer.py}, documenting the HMAC key
rotation risk: a payload serialised under a rotated-out key cannot be rejected
as forged, but it can be rejected as expired when the loading site validates
the embedded timestamp.

\paragraph{Phase~2 Baseline.}

\begin{center}
\begin{tabular}{ll}
\toprule
Metric & Value \\
\midrule
Tests passed & 3,925 \\
Tests skipped & 120 \\
Tests failed & 0 \\
Branch coverage & 98.26\,\% \\
Python & 3.13.7 \\
pytest & 7.4.4 \\
Natural policy tests & 61 \\
Wall time & $\approx$\,430\,s \\
\bottomrule
\end{tabular}
\end{center}

\paragraph{Post-Upgrade Baseline (pytest 8.4.2, 2026-05-13).}

\begin{center}
\begin{tabular}{ll}
\toprule
Metric & Value \\
\midrule
Tests passed & 4,002 \\
Tests skipped & 81 \\
Tests failed & 0 \\
Items collected & 4,082 \\
Branch coverage & 98.26\,\% \\
Python & 3.13.7 \\
pytest & 8.4.2 \\
Wall time & 437.91\,s \\
\bottomrule
\end{tabular}
\end{center}

The pytest upgrade from 7.4.4 to 8.4.2 changed skip semantics for several
\code{anyio}-backed fixtures (anyio 4.12.1, asyncio-mode \texttt{AUTO});
39 previously-skipped tests now execute and pass, raising the passing count
from 3,925 to 4,002 with no new failures.

% ───────────────────────────────────────────────────────────────────────────────
\subsection*{Phase~3 (Pillar~1 IR Compiler) --- v1.2: Strict IR Compiler and Decompiler}
% ───────────────────────────────────────────────────────────────────────────────

\paragraph{Motivation.}
Phase~2 (\fname{natural\_policy/}) established a text-to-policy pipeline but
required an active LLM client and produced output tied to the
\code{NaturalPolicySchema} intermediate representation.  For deployments
already receiving structured JSON from the LLM via Structured Outputs or
JSON Mode, the full Phase~2 pipeline is heavier than necessary.  More
importantly, Phase~2's compiler lacks field-to-field comparison support,
recursive rule nesting, and \code{IN}/\code{NOT\_IN} membership operators.
Pillar~1 addresses all three, with no LLM dependency of its own.

\paragraph{Design Principles.}

\begin{enumerate}
\item \textbf{Standalone.}  \code{PolicyCompiler} has zero imports from any
  LLM library or from \fname{natural\_policy/}.  It operates on a
  fully-validated \code{PolicyIR} object.
\item \textbf{No LLM on the hot path.}  The compiler, IR models, and
  \code{Decompiler} have no dependency on any LLM client library.  The LLM's
  only role is to produce the \code{PolicyIR} JSON via Structured Outputs;
  the compiler takes it from there.
\item \textbf{Richer IR.}  \code{PolicyIR} supports field-to-field
  comparisons (\code{FieldReference} RHS), recursive rule nesting,
  \code{IN}/\code{NOT\_IN} membership operators, and eight binary comparison
  operators.
\item \textbf{CISO auditability.}  Every compiled policy can be decompiled
  back to structured English via \code{Decompiler} --- a timestamped report
  for CISO review before a dynamically compiled policy is deployed.
\end{enumerate}

\paragraph{Module Architecture.}

\[
  \underbrace{\text{LLM structured output}}_{\text{JSON}}
  \;\xrightarrow{\text{Pydantic validate}}\;
  \code{PolicyIR}
  \;\xrightarrow{\code{PolicyCompiler}}\;
  \code{list[ConstraintExpr]}
  \;\xrightarrow{\code{Decompiler}}\;
  \text{CISO report}
\]

\code{PolicyIR.model\_json\_schema()} is the \code{response\_format} passed
to the LLM.  Any malformed LLM output is rejected by Pydantic~v2 before
\code{PolicyCompiler} is ever invoked.

\paragraph{Compile-Time Error Classes.}
Seven categories of policy error are detected before any Z3 formula is
constructed:
\begin{enumerate}
\item Unknown field referenced (LLM hallucination).
\item Scalar type incompatible with field Z3 sort (\code{String} field vs.\
  integer literal).
\item Field-to-field sort mismatch (\code{Real}-sorted vs.\ \code{Bool}-sorted).
\item Ordering operator (\code{>}, \code{<}, \code{>=}, \code{<=}) on a
  \code{Bool}-sorted field.
\item Membership operator (\code{IN}/\code{NOT\_IN}) with a non-list RHS.
\item Empty membership set (vacuously unsatisfiable).
\item \code{bool} literal as RHS of an ordering operator.
\end{enumerate}
All seven raise \code{PolicyCompilationError} with full context.  No partial
result is ever returned --- fail-closed is structural, not conventional.

\paragraph{Public Surface.}
Nine symbols are exported via both \code{compiler.\_\_all\_\_} and
\code{pramanix.\_\_init\_\_}:
\code{Condition}, \code{Decompiler}, \code{FieldReference},
\code{FieldSource}, \code{Logic}, \code{Operator}, \code{PolicyCompiler},
\code{PolicyIR}, \code{Rule}.

\subsection*{Phase~12: Metric-Registry Bug and Coverage Completeness}

Phase~12 closed a silent production defect in the Prometheus instrumentation
and plugged all remaining testable coverage gaps.

\paragraph{Bug~3 --- Prometheus Counter Double-Registration in
  \code{worker.py}.}
\code{\_warmup\_worker()} created a \code{prometheus\_client.Counter} object
\emph{inside} the \code{except~Exception} handler.  In \code{async-thread}
mode every thread worker shares the same module namespace; the second
invocation of \code{\_warmup\_worker()} raised \code{ValueError: Duplicated
timeseries in CollectorRegistry}.  That exception was silently swallowed by
the outer \code{except~Exception:~pass} clause, permanently disabling warmup
failure counting for every call after the first.  The defect was latent in
single-process mode but would corrupt metrics in any multi-worker deployment.

\textbf{Fix.}  A module-level \code{\_WORKER\_WARMUP\_FAILURE\_COUNTER}
variable is now initialised once at import time, following the identical
pattern already used by \code{crypto.py} (lines~55--71).  A
\code{try/except~ValueError} block attempts to register the counter; if the
name already exists in \code{prometheus\_client.REGISTRY} the reference is
recovered via \code{REGISTRY.\_names\_to\_collectors}.  The
\code{\_warmup\_worker()} exception handler now simply calls
\code{\_WORKER\_WARMUP\_FAILURE\_COUNTER.inc()}, guarded by a
\code{is~not~None} check and an inner \code{try/except~Exception} so that a
missing Prometheus installation never crashes the worker.

\paragraph{Coverage Gaps Closed.}

\begin{itemize}

\item \textbf{Haystack adapter async error paths}
  (\fname{integrations/haystack.py} lines 176, 181--187, 208--209).\quad
  Six new tests in \fname{tests/integration/test\_haystack\_adapter.py}
  (class \code{TestRunAsync}) exercise every branch of the async error path.
  The \code{\_AsyncBrokenGuard} helper class provides a real
  \code{async~def~verify\_async()} implementation that raises
  \code{RuntimeError}, satisfying the no-mock constraint while covering
  both the \code{block\_on\_error=True} and \code{block\_on\_error=False}
  branches.  The \code{messages}/\code{blocked\_messages} assignment branches
  (lines~208--209) are covered by tests that supply and omit the optional
  \code{messages} parameter.

\item \textbf{LangChain execute-map dispatch branch}
  (\fname{integrations/langchain.py} line~178).\quad
  One new test in \fname{tests/integration/test\_langchain\_tool.py}
  (\code{test\_wrap\_tools\_uses\_execute\_map\_when\_tool\_name\_matches})
  passes an \code{execute\_map} dictionary whose key matches the tool name.
  It asserts that the wrapped tool's \code{\_pramanix\_execute} attribute is
  exactly the function supplied via the map rather than the default
  \code{tool.run} callable.

\item \textbf{LlamaIndex async state provider}
  (\fname{integrations/llamaindex.py} line~484).\quad
  One new test in \fname{tests/integration/test\_integration\_coverage.py}
  (\code{test\_acall\_with\_async\_state\_provider}) passes an
  \code{async~def} coroutine as the \code{state\_provider} argument to
  \code{PramanixQueryEngineTool}.  The test exercises the
  \code{if asyncio.iscoroutine(result): result = await result} branch
  inside \code{\_get\_state()}.

\end{itemize}

\paragraph{Remaining Structurally Dead Lines.}  The following lines cannot be
covered without modifying the production environment and are documented as
intentional:

\begin{itemize}
\item \fname{worker.py}~242--246: the ppid watchdog tight loop (only executes
  inside a live subprocess spawned by \code{ProcessPoolExecutor}).
\item \fname{worker.py}~85--95: the \code{except~ValueError} registry
  recovery path (only reachable when \code{\_warmup\_worker} is called a
  second time inside the same Python process during multi-worker startup
  on a platform where Prometheus is installed — rare in CI).
\item \fname{archiver.py}~161--164, 196--199: \code{cryptography} package
  absent and OS-level \code{fsync}/\code{replace} failures.
\item \fname{transpiler.py}~346: structurally unreachable \code{continue}
  following an exhaustive \code{if/elif/else}.
\item All integration import-error placeholder blocks: require uninstalling
  the corresponding framework package mid-test.
\item \fname{cli.py}~1033--1054: Linux/musl platform detection (Windows CI).
\item \fname{gemini.py}~199--221: \code{httpx} transport-level error
  injection.
\end{itemize}

\paragraph{Coverage Impact.}  Phase~12 adds 7 new tests (3,923 total) and
improves branch coverage from 98.22\,\% to \textbf{98.26\,\%}.  The 98\,\%
gate remains satisfied with margin.

\subsection*{Phase~11: Test-Isolation Hardening}

A post-release audit revealed two latent test bugs that only surfaced when the
full 3,916-test suite ran in its default ordering.  Both were genuine
production-quality defects in the test infrastructure, not in Pramanix itself.

\paragraph{Bug~1 --- Missing \code{\_dlq\_pending} Initialisation.}
\code{TestKafkaConsumerRealPaths.\_make\_consumer()} used
\code{PramanixKafkaConsumer.\_\_new\_\_} to bypass the real Kafka bootstrap,
but omitted the \code{\_dlq\_pending~= 0} and \code{\_dlq\_flush\_interval~= 100}
attribute initialisations that \code{\_\_init\_\_} normally performs.  When a
guard exception caused \code{safe\_poll} to call \code{\_dead\_letter()},
\code{self.\_dlq\_pending~+= 1} raised \code{AttributeError}, masking the
underlying guard error and causing the generator to yield nothing.  Fix: both
attributes are now initialised explicitly in \code{\_make\_consumer()}.

\paragraph{Bug~2 --- Structlog State Leaking Across Test Files.}
\code{TestLoggingIsolation.\_capture\_log\_output()} in
\fname{test\_guard\_dark\_paths.py} called \code{structlog.configure(\ldots,
logger\_factory=PrintLoggerFactory(buf))} without restoring the prior
configuration.  When \code{test\_guard\_dark\_paths.py} ran \emph{before}
\fname{test\_interceptors\_real.py} in the suite ordering, the leaked
\code{PrintLoggerFactory} remained active.  Structlog's \code{add\_logger\_name}
processor subsequently called \code{logger.name} on a \code{PrintLogger}
(which lacks a \code{.name} attribute), raising \code{AttributeError} inside
every \code{guard.verify()} call made by the Kafka consumer tests.  The
exception was swallowed by \code{safe\_poll}'s \code{except~Exception} clause,
dead-lettering all messages and causing the generators to yield nothing.
Fix: wrapped the \code{structlog.configure()} call in a \code{try/finally}
block that captures the old config via \code{structlog.get\_config()} and
restores it unconditionally after the test, making each invocation
of \code{\_capture\_log\_output()} side-effect-free.

\paragraph{Coverage Impact.}  With both bugs fixed, line coverage increased
from 98.16\,\% to \textbf{98.22\,\%} (0\,failures, 98\,\% gate: passing).

\subsection*{Phase~10: Production Hardening}

Phase~10 addressed three root-cause gaps identified during the final pre-PyPI
audit: two open Known Limitations, three test failures, and a coverage gate
that was 0.31\,\% short.

\paragraph{AES-256-GCM Archive Encryption.}
\code{EncryptedArchiveWriter} was added to \fname{audit/archiver.py}.  It
implements the pluggable \code{archive\_writer: Callable[[Path, bytes], None]}
protocol defined in Phase~8.  Each archive segment is encrypted with a fresh
96-bit nonce generated by \code{secrets.token\_bytes(12)} (NIST~SP~800-38D
\S8.2 compliant).  The GCM authentication tag is embedded in the ciphertext;
any byte-level tampering raises \code{InvalidTag} before plaintext is exposed.
Writes are atomic: \code{tempfile.mkstemp()} + \code{os.fdopen()} +
\code{os.fsync()} + \code{os.replace()} ensures no partial segment is ever
visible on disk.  The class is exported from \code{pramanix.audit} and
documented in \fname{docs/PUBLIC\_API.md} as Tier~B (stable).  Known
Limitation~5 (Merkle Archive Plaintext) is resolved.

\paragraph{Configurable Injection-Score Field Name.}
\code{injection\_confidence\_score()} in \fname{translator/\_sanitise.py} now
accepts \code{*, amount\_field: str = "amount"} as a keyword-only parameter.
Non-financial callers pass the field name that carries a monetary amount in
their schema; passing \code{amount\_field=""} disables the sub-penny signal
entirely.  The change is fully backward-compatible: callers that do not pass
\code{amount\_field} receive the same behaviour as before.  Known
Limitation~9 (Injection Scorer Sub-Penny Gap) is resolved.

\paragraph{Three Test Failures.}
All three failures were root-cause-diagnosed and fixed with surgical changes:
(1)~a missing \code{\_SKLEARN\_AVAILABLE} guard in \code{CalibratedScorer.load()};
(2)~a monkey-patch-vs-call-site binding error fixed with
\code{unittest.mock.patch()} at the correct import path;
(3)~\code{GeminiTranslator.\_\_init\_\_}'s except clause broadened to
\code{(ImportError, AttributeError)}.  No test intent was changed.

\paragraph{Coverage Gate.}
Eleven new real-integration tests in
\fname{tests/unit/test\_archiver\_coverage\_gaps.py} exercising the full
\code{EncryptedArchiveWriter} lifecycle --- key generation, round-trip
encrypt/decrypt, atomic write, wrong-key \code{InvalidTag}, fresh-nonce-per-call
--- lifted branch coverage from 97.69\,\% to 98.16\,\%, crossing the gate.

%% ============================================================================
\chapter{Testing Strategy}
%% ============================================================================

\section{Test Architecture Philosophy}

The test suite is not a collection of happy-path unit tests. It is an
adversarial testing infrastructure, documented in \fname{reality.md}
\cfrsec{3.1}.

\section{Adversarial Tests}

\textbf{Z3 Context Isolation}
(\fname{tests/adversarial/test\_z3\_context\_isolation.py}): Uses
\code{threading.Barrier} to force 10 concurrent simultaneous solver executions.
Demonstrates empirically that Z3 context state does not bleed between threads
under the tested conditions (10 concurrent solvers, single process).  This is
adversarial evidence of isolation, not a mathematical proof of impossibility;
a formal proof would require verification of Z3's internal C implementation.
CVE-prevention level adversarial testing.

\textbf{Prompt Injection}
(\fname{tests/adversarial/test\_prompt\_injection.py}): Verifies resistance to
null bytes, unicode full-width digits, massive integers, negative bounds, and
resource exhaustion using in-process stub translators. Isolates the Z3 +
governance layer from LLM parsing --- correct layered unit-test design.

\section{Property-Based Tests (Hypothesis)}

\code{hypothesis} is in the dev dependencies (\fname{pyproject.toml}:
\code{hypothesis = "\^{}6.100"}). Property-based testing verifies that the
DSL produces correct constraints for arbitrary input values rather than just
hand-picked examples.

Known gap (\fname{reality.md} \cfrsec{3.1}): The \fname{\_sanitise.py}
injection scorer has complex additive float-math heuristics. No Hypothesis
property tests verify that score combinations are monotonically bounded within
$[0.0, 1.0]$ across the full input space.

\section{Real Infrastructure Integration Tests}

\fname{tests/integration/} --- 26 files hitting real containers: Kafka via
Redpanda, Redis, Postgres, LocalStack, Azure KeyVault, HashiCorp Vault. Live
LLM adapters tested: Gemini, Cohere. Container availability is detected at
test session start (\fname{tests/integration/conftest.py}:34--45). The 81
skips in the test suite are primarily these infrastructure-dependent tests
skipping when Docker is unavailable. This is correct behavior, not a testing
gap.

\section{Concurrent Load Tests}

\fname{tests/unit/test\_production\_gaps.py} --- 14 tests including a
50-coroutine concurrent async size-check test with oversized payloads,
unserializable payloads, and mixed batches. Timing-pad distribution tests:
p5 latency assertions ($\geq 90$\,\% of budget) for both ALLOW and BLOCK;
symmetry ratio check ($\leq 1.30$) to detect asymmetric padding.

\section{Coverage Gate}

\fname{pyproject.toml}: \code{fail\_under = 98}.  Current status after
Phase~11: \textbf{98.22\,\% --- gate passing}.  Branch coverage is tracked
alongside line coverage (\code{--cov-branch}).

Residual uncovered lines fall into structurally-dead categories that
cannot be reached without removing installed packages or injecting OS-level
failures:
\begin{itemize}
  \item \fname{audit/archiver.py}: lines~161--164 (the
        \code{ConfigurationError} branch that fires when the
        \code{cryptography} package is absent) and lines~196--199 (the
        \code{os.write()} failure path inside the atomic write helper).
  \item \fname{transpiler.py}: line~346 (a \code{continue} in the final
        enum-promotions loop that is structurally unreachable with any
        valid input).
  \item Optional-dependency import-error paths in all six integration
        adapters: \fname{langchain.py}, \fname{haystack.py},
        \fname{llamaindex.py}, \fname{crewai.py},
        \fname{kafka.py}, \fname{fastapi.py}.
        These paths execute only when the optional framework is \emph{not}
        installed; covering them would require uninstalling those packages
        for a dedicated test run.
  \item \fname{translator/gemini.py}: lines~199--221 (\code{httpx} transport
        error wrapping and \code{AsyncRetrying} exception re-raise path).
        Requires a network-level timeout injection that cannot be simulated
        without the real Gemini endpoint or a live httpx mock server.
  \item \fname{worker.py}: lines~242--246, 338 (\code{WorkerPool} crash-
        restart path and an orphaned \code{asyncio.Task} finaliser).
        Triggered only on OS-level thread scheduling failures.
\end{itemize}

\section{All Test Failures --- Resolved (Phases~10 and~11)}

All five diagnosed test failures have been fixed.  The test suite reports 0
failures.  Phase~10 resolved the three v1.0.0 originals; Phase~11 resolved two
latent ordering-dependent bugs discovered during the post-release audit.

\textbf{Failure~1 ---
\code{test\_sklearn\_absent\_raises\_configuration\_error}} (\textit{fixed}):
\code{CalibratedScorer.load()} classmethod lacked the sklearn availability
check present in \code{\_\_init\_\_()}.  Fix: added the same
\code{\_SKLEARN\_AVAILABLE} guard at the start of \code{load()}, raising
\code{ConfigurationError} before any attribute access.

\textbf{Failure~2 ---
\code{test\_canonical\_bytes\_exception\_falls\_back\_to\_stdlib\_json}}
(\textit{fixed}): The test monkey-patched \code{\_dec\_mod.\_canonical\_bytes}
after import, which had no effect on the already-resolved binding inside
\code{\_compute\_hash()}.  Fix: replaced the monkey-patch with
\code{unittest.mock.patch("pramanix.decision.\_canonical\_bytes",
side\_effect=RuntimeError(...))} which patches at the actual call site.

\textbf{Failure~3 --- \code{test\_gemini\_prefix}} (\textit{fixed}):
\code{GeminiTranslator.\_\_init\_\_} only caught \code{ImportError}, but the
\code{google} namespace package (installed via \code{google.auth}) caused an
\code{AttributeError} inside \code{google-generativeai}.  Fix: changed the
catch clause to \code{(ImportError, AttributeError)} and changed
\code{\_GEMINI\_AVAILABLE} detection from \code{find\_spec()} to an actual
try-import, making the availability check consistent with the runtime
behaviour.

\textbf{Phase~11 --- Failure~4 ---
\code{test\_safe\_poll\_allowed\_message\_yielded}} (\textit{fixed}):
See Phase~11 subsection above (Bug~1, missing \code{\_dlq\\_pending}
initialisation in \code{\\_make\\_consumer()}).

\textbf{Phase~11 --- Failure~5 ---
\code{test\_safe\_poll\_guard\_exception\_dead\_letters\_and\_continues}}
(\textit{fixed}): See Phase~11 subsection above (Bug~2, structlog
\code{PrintLoggerFactory} state leak contaminating Kafka consumer tests).

%% ============================================================================
\chapter{Research and Benchmarks}
%% ============================================================================

\section{Latency Benchmarks}

\textbf{Benchmark environment}: Windows~11 (Build~10.0.26200), AMD64,
Intel Core 14P/20-logical-thread CPU at 2300\,MHz base clock, 15.6\,GB RAM
(consumer laptop hardware --- not a server or cloud instance).  Single
process, single thread, \code{sync} execution mode, Python~3.13.7,
Z3~4.16.0, \code{orjson} installed, \code{prometheus-client} installed.
The 1M stability test was collected under Pramanix~v0.8.0; the current
release is v1.0.0.  All figures are wall-clock elapsed time, not CPU time.
Results are not portable across hardware without re-running
\fname{benchmarks/latency\_benchmark.py}.

Results from \fname{benchmarks/results/latency\_results.json} --- API mode,
2,000 samples, 5-invariant policy, \code{sync} execution mode:

\begin{table}[h]
\centering
\begin{tabular}{@{} l r r @{}}
\toprule
\textbf{Metric} & \textbf{Measured} & \textbf{Target} \\
\midrule
P50  & 5.235\,ms & 5.0\,ms \\
P95  & 6.361\,ms & 10.0\,ms \\
P99  & 7.109\,ms & 15.0\,ms \\
Mean & 5.336\,ms & --- \\
\bottomrule
\end{tabular}
\caption{Latency benchmark results (source:
  \fname{benchmarks/results/latency\_results.json})}
\label{tab:latency}
\end{table}

P50 narrowly misses its target by 0.235\,ms. P95 and P99 are well within
targets. The \code{passed: false} in the JSON file reflects the strict P50
gate.

\section{1M Decision Stability Test}

Results from \fname{benchmarks/results/1m\_audit\_summary.json} ---
1,000,000 decisions, 500 warmup, 1-invariant policy, \code{sync} mode,
single-threaded:

\begin{table}[h]
\centering
\begin{tabular}{@{} l r @{}}
\toprule
\textbf{Metric} & \textbf{Value} \\
\midrule
Total decisions    & 1,000,000 \\
Total elapsed      & 12,298.479\,s \\
Average RPS        & 81.3 \\
P50 latency        & 11.283\,ms \\
P95 latency        & 20.145\,ms \\
P99 latency        & 30.538\,ms \\
P99.9              & 153.848\,ms \\
P99.99             & 270.578\,ms \\
Max                & 1,565.746\,ms \\
Baseline RSS       & 57.617\,MiB \\
Final RSS          & 60.422\,MiB \\
Peak RSS           & 80.395\,MiB \\
Memory growth      & 2.805\,MiB over 1\,M decisions \\
GC gen0 delta      & $+$6 collections \\
GC gen1/gen2 delta & 0 \\
\bottomrule
\end{tabular}
\caption{1M decision stability test (source:
  \fname{benchmarks/results/1m\_audit\_summary.json})}
\label{tab:stability}
\end{table}

The memory profile is critical evidence. 2.805\,MiB growth across 1,000,000
decisions is not a leak. The RSS oscillations visible in
\fname{1m\_audit\_full.log} (spiking between 57\,MiB and 74\,MiB throughout
the run) are Z3's native heap behavior --- the solver allocates and deallocates
Z3 context objects in its internal C heap, and Python's process RSS reflects
those allocations with some lag. The final RSS is only 2.8\,MiB above
baseline. This is bounded, stable behavior.

\textbf{P50 discrepancy between benchmark tables.}  The 2,000-sample latency
benchmark (Table~\ref{tab:latency}) records P50 of 5.235\,ms on a
5-invariant policy.  The 1M stability benchmark (Table~\ref{tab:stability})
records P50 of 11.283\,ms on a 1-invariant policy --- a simpler policy that
runs 2.2$\times$ slower at the median.  The most likely explanation: the 1M
run measures steady-state performance under sustained load, where GC pressure,
Z3 internal heap fragmentation, and OS scheduling jitter accumulate over
$\sim$3.4 hours.  The 2,000-sample run completes in approximately ten seconds
and captures warm-cache, low-GC conditions.  Both benchmarks ran on the same
hardware in the same \code{sync} mode.  The figures serve different purposes:
the latency benchmark characterises request-path overhead; the 1M benchmark
characterises long-run stability.  Neither figure should be used as a latency
SLA without hardware-matched reproduction.

\section{100M Decision Architecture}

\fname{benchmarks/100m\_orchestrator\_fast.py} and
\fname{benchmarks/100m\_worker\_fast.py} document a zero-IPC production
architecture:

\begin{lstlisting}[style=plain,caption={100M decision target architecture}]
18 x multiprocessing.Process (spawn mode)
Each process: owns sync Guard + Z3 in one address space
Decision loop: cycles through 10,000 pre-built payloads (5 MiB per worker)
Target: 100M decisions per domain x 5 domains = 500M decisions total
\end{lstlisting}

\begin{table}[h]
\centering
\begin{tabular}{@{} p{6.5cm} r r @{}}
\toprule
\textbf{Architecture} & \textbf{RPS/worker} & \textbf{IPC crossings} \\
\midrule
Old: asyncio $\to$ ThreadPool $\to$ Guard process $\to$ Z3 process
  & 22--27 & 2 per decision \\
New: 18 process $\times$ sync Guard
  & 80--120 & 0 per decision \\
\bottomrule
\end{tabular}
\caption{Architecture throughput comparison (source:
  \fname{benchmarks/100m\_worker\_fast.py})}
\end{table}

The 3.5--5$\times$ throughput improvement comes entirely from eliminating IPC
crossings. 18 workers targeting 80--120~RPS each yields 1,440--2,160
aggregate~RPS.

%% ============================================================================
\chapter{Use Cases and Target Industries}
%% ============================================================================

\section{Financial Services}

\textbf{Banking Transfers}: \fname{examples/banking\_transfer.py} is the
canonical reference. A \code{BankingPolicy} with three invariants
(\code{non\_negative\_balance}, \code{within\_daily\_limit},
\code{account\_not\_frozen}) demonstrates all six decision outcomes. Exact
arithmetic via \code{Decimal} prevents the IEEE~754 rounding errors that
would cause $\mathit{balance} - \mathit{amount} \geq 0$ to fail or pass
incorrectly.

\textbf{High-Frequency Trading --- Wash Sale Detection}:
\fname{examples/hft\_wash\_trade.py}. The \code{WashSaleDetection} primitive
in \fname{primitives/fintech.py} directly maps to IRC~\S\,1091.

\textbf{Anti-Structuring / BSA-AML}: \code{AntiStructuring} in
\fname{primitives/fintech.py}. Maps to BSA/AML 31~CFR~\S\,1020.320(a)(2). An
AI agent structuring deposits to avoid BSA reporting thresholds is a criminal
violation. Z3-verified anti-structuring invariants enforce a threshold check consistent
with the reporting conditions in 31~CFR~\S\,1020.320(a)(2).  Whether that
check constitutes legal compliance in a given deployment depends on the
broader compliance program and regulatory examiner judgment; this SDK alone
does not confer legal compliance.

\textbf{Collateral Management}: \code{CollateralHaircut},
\code{MarginRequirement}, \code{MaxDrawdown}, \code{VelocityCheck} ---
quantitative invariants that Z3 handles exactly.

\section{Healthcare}

\textbf{PHI Access Control}: \fname{examples/healthcare\_phi\_access.py} and
\fname{examples/healthcare\_rbac.py}. \code{PHILeastPrivilege},
\code{BreakGlassAuth}, \code{ConsentActive} enforce access controls
consistent with the HIPAA minimum-necessary principle at the action boundary.
Every PHI access produces a signed, Merkle-anchored \code{Decision} with a
\code{decision\_hash} --- suitable for incorporation into a HIPAA audit trail
programme.  See the \textit{Regulatory Mapping} chapter for scope and
limitations; these records alone do not satisfy all HIPAA audit logging
obligations.

\textbf{Dosage Safety}: \code{DosageGradientCheck},
\code{PediatricDoseBound} --- invariants over prescription amounts. An AI
clinical decision support system that recommends dosages cannot exceed
pediatric weight-based bounds. Z3 verifies this before the recommendation
reaches the prescriber.

\section{Cloud Infrastructure}

\textbf{Blast Radius Control}: \fname{examples/infra\_blast\_radius.py}.
\code{BlastRadiusCheck} in \fname{primitives/infra.py}. An AI agent scaling
Kubernetes pods cannot set replicas above a hard cap without a
\code{ProdDeployApproval} invariant passing. \code{ReplicaBudget},
\code{CPUMemoryGuard}, \code{WithinCPUBudget}, \code{WithinMemoryBudget} ---
all computable from current cluster state and expressible as Z3 Real
arithmetic.

\textbf{SRE Automation}: Infrastructure changes proposed by AI SRE systems
are Z3-verified before applying. A change that would violate
\code{MinReplicas} during a traffic spike is blocked with a full explanation.

\section{Multi-Agent Systems}

\textbf{AutoGen / CrewAI}: \fname{examples/autogen\_multi\_agent.py}. The
\code{PramanixCrewAITool} and AutoGen adapter ensure that every tool
invocation in a multi-agent system --- not just the first --- is verified. An
agent in a chain cannot bypass a policy by delegating to a sub-agent.

\textbf{LangChain / LlamaIndex}:
\fname{examples/langchain\_banking\_agent.py},
\fname{examples/llamaindex\_rag\_guard.py}. A retrieval-augmented agent
querying a document database and proposing an action based on retrieved
content cannot use that content to bypass the policy.

%% ============================================================================
\chapter{Known Limitations and Future Roadmap}
%% ============================================================================

\section{Verified Limitations (from \fname{docs/KNOWN\_GAPS.md})}

\begin{enumerate}
  \item \textbf{AGPL-3.0 License} (\fname{pyproject.toml}:8). Most enterprise
        legal teams in regulated industries do not approve AGPL as a Python
        library dependency without explicit exception review, due to the
        network-service copyleft obligation. The entire target market ---
        financial services, healthcare, cloud providers --- defaults to
        Apache~2.0, MIT, or commercial licenses. This is documented in
        \fname{reality.md} \cfrsec{4.2} as the first adoption blocker.

  \item \textbf{Redis Circuit Breaker TOCTOU Race} --- \textit{resolved}.
        \code{RedisDistributedBackend.set\_state}
        (\fname{circuit\_breaker.py}:771) uses \code{WATCH} +
        \code{MULTI}/\code{EXEC} optimistic locking to make the entire
        read-modify-write atomic.  If another writer modifies the key between
        \code{WATCH} and \code{EXEC}, the transaction aborts and retries up
        to \code{\_MERGE\_MAX\_RETRIES=3} times (\fname{circuit\_breaker.py}:769).
        No Lua scripting required.

  \item \textbf{\code{\_tree\_repr} Polynomial Gap} --- \textit{resolved}.
        \fname{transpiler.py} contains explicit \code{match}/\code{case}
        clauses for both \code{\_PowOp} and \code{\_ModOp} (lines~762--765).
        Structural fingerprints for polynomial and modulo policies are correct.

  \item \textbf{Logging Split-Brain} (\fname{reality.md} \cfrsec{4.1}).
        \fname{guard\_config.py} installs structlog with secrets redaction.
        But \fname{worker.py}, \fname{solver.py}, \fname{decision.py} and most
        other modules use \code{logging.getLogger(\_\_name\_\_)} (stdlib),
        bypassing structlog redaction entirely. Sensitive policy values can
        reach log sinks unredacted.

  \item \textbf{Merkle Archive Encryption} (\fname{audit/archiver.py}).
        \textit{[Resolved — Phase~10]}. AES-256-GCM at-rest encryption is
        provided by \code{EncryptedArchiveWriter}.  Pass it as
        \code{archive\_writer=EncryptedArchiveWriter(key)} at construction.
        The archiver emits a \code{WARNING} only when no writer is supplied.
        SOC~2 CC6.6, PCI~DSS Req~10, and HIPAA \S164.312(a)(2)(iv) are
        now addressable at the SDK layer.

  \item \textbf{HMAC-Only JWT} (\fname{identity/linker.py}). HMAC-SHA256
        symmetric. Every service in a microservice mesh must share the secret
        key. RS256 or ES256 required for enterprise deployments.

  \item \textbf{Execution Token Replay After Restart}
        (\fname{execution\_token.py}). The consumed-set is in-memory only. A
        process restart clears it. \code{RedisExecutionTokenVerifier} closes
        this for Redis deployments.

  \item \textbf{No NLP Validators}. No PII redaction, toxicity scoring, or
        free-text classification. Pramanix is built exclusively for structured
        intent verification. Pramanix~$+$ Guardrails~AI are complementary,
        not alternatives.

  \item \textbf{Injection Scorer Field Name}
        (\fname{\_sanitise.py}). \textit{[Resolved — Phase~10]}.
        \code{injection\_confidence\_score} now accepts a keyword-only
        \code{amount\_field: str = "amount"} parameter.  Non-financial
        callers pass their field name directly; passing
        \code{amount\_field=""} disables the sub-penny signal entirely
        without needing to set \code{sub\_penny\_threshold=Decimal("0")}
        as a workaround.  The API is backward-compatible: existing callers
        need no changes.

  \item \textbf{Stub Integrations} (\fname{integrations/}). DSPy, Haystack,
        Semantic Kernel, and Pydantic AI are stubs
        (\fname{docs/KNOWN\_GAPS.md} \cfrsec{8}).
\end{enumerate}

\section{Future Roadmap}

\textbf{Immediate (Pre-PyPI):}
\begin{itemize}
  \item License change: AGPL-3.0 $\to$ Apache~2.0
  \item Redis circuit breaker TOCTOU already resolved via \code{WATCH}/\code{MULTI}/\code{EXEC} (\fname{circuit\_breaker.py}:771)
  \item \code{\_tree\_repr} polynomial gap already resolved (\fname{transpiler.py}:762)
  \item \textit{[Done — Phase~10]} Fix 3 failing tests, restore 98\,\% coverage gate --- 3,916 passed, 81 skipped, 0 failed; 98.16\,\% coverage
  \item \textit{[Done — Phase~10]} Merkle archive pluggable AES-256-GCM writer (\code{EncryptedArchiveWriter})
  \item \textit{[Done — Phase~10]} Configurable \code{amount\_field} in injection confidence scorer
  \item \textit{[Done — Phase~11]} Fix 2 latent test-isolation bugs; coverage: 98.16\,\% $\to$ 98.22\,\%
\end{itemize}

\textbf{Short-Term (v1.1):}
\begin{itemize}
  \item License change: AGPL-3.0 $\to$ Apache~2.0 (unchanged blocker)
  \item Logging split-brain fix: route stdlib logging through
        \code{structlog.stdlib.ProcessorFormatter}
  \item RS256/ES256 JWT support
  \item Hypothesis property tests for injection scorer
  \item Close optional-dependency coverage gaps in
        \fname{integrations/langchain.py},
        \fname{integrations/haystack.py},
        \fname{integrations/llamaindex.py},
        \fname{integrations/crewai.py},
        \fname{interceptors/kafka.py} (currently 86--90\,\%)
\end{itemize}

\textbf{Medium-Term (v1.2):}
\begin{itemize}
  \item Complete DSPy, Semantic Kernel, Pydantic AI integrations
  \item Z3 formula caching across Guard instances with identical policies
  \item Declarative policy YAML front-end (without removing Python DSL)
\end{itemize}

%% ============================================================================
\chapter{Competitive Deep-Dive}
%% ============================================================================

\section{Open Policy Agent (OPA)}

\textbf{What OPA does}: Policy-as-code with Rego, a Datalog-inspired
declarative language. Evaluates structured JSON data against policy rules.
Returns allow/deny. Used widely for Kubernetes admission control and API
authorization.

\textbf{Where OPA falls short for Pramanix's use case}: OPA has built-in
arithmetic operators and can express numeric comparisons in Rego.  The
constraint $\mathit{balance} - \mathit{amount} \geq \mathit{min\_reserve}$
is expressible in Rego.  The substantive limitations are:
\begin{enumerate}
  \item \textbf{No formal satisfiability guarantee.}  A Rego rule returning
        \code{false} means ``no derivation succeeded'' --- not ``this
        arithmetic constraint is violated by the inputs.''  There is no
        underlying SMT completeness proof.
  \item \textbf{No unsat-core attribution.}  OPA has no equivalent to Z3's
        \code{unsat\_core()}.  Identifying which specific rule caused a
        denial requires manual rule-set introspection or explicit trace
        annotations, not a built-in mechanism.
  \item \textbf{Domain mismatch.}  OPA is optimized for authorization policy
        (who may do what).  Quantitative safety verification --- do the
        arithmetic relationships between fields satisfy all constraints
        simultaneously? --- is a different problem class.
\end{enumerate}

\textbf{Correct use}: OPA for authorization (\textit{can user X perform
action~Y?}). Pramanix for mathematical safety (\textit{is action~Y logically
consistent with all invariants?}). Use both.

\section{Guardrails AI}

\textbf{What Guardrails AI does}: Provides \code{Validator} objects that run
checks on LLM input/output strings. Hundreds of pre-built validators: PII
detection, toxicity, content safety, format checks.

\textbf{Where Guardrails AI falls short}: Guardrails AI is a text-processing
library. Its validators are probabilistic functions. \code{(E(balance) -
E(amount) >= 0)} cannot be expressed as a Guardrails validator without
reimplementing Z3 arithmetic.

\textbf{Market reality}: The two systems are complementary. Guardrails~AI for
content safety on LLM I/O. Pramanix for formal safety on structured actions.

\section{NVIDIA NeMo Guardrails}

NeMo Guardrails is a dialogue management framework using Colang-based
conversation flow control.  Pramanix is transactional action verification.
They operate in different problem domains.  NeMo manages conversation-level
guardrails (topic steering, topical restrictions); Pramanix verifies structured
action intent against arithmetic invariants at tool call boundaries.  A
deployment can use both without conflict.  No independent benchmarking
comparing the two systems on a common task has been conducted; comparative
claims about relative architectural merit without such data are unsupported.

\section{Direct AI Safety Tools}

Several tools address LLM output and agent safety that are not covered in the
sections above:

\textbf{Protect AI LLM Guard} (\url{https://llm-guard.com}): Input/output
scanning for PII, prompt injection, toxicity, and harmful content.  Text-based
scanning using NLP classifiers; does not perform Z3 arithmetic verification
of structured actions.  Complementary to Pramanix.

\textbf{LangSmith Guardrails}: Evaluation and tracing framework within the
LangChain ecosystem.  Focuses on observability and evaluation; does not
provide formal invariant verification.

\textbf{Rebuff} (\url{https://github.com/protectai/rebuff}): Prompt injection
detection using heuristics and LLM-based classifiers.  Complementary to
Pramanix's injection defence stack.

\textbf{Guardrails Hub}: Community-contributed validators for Guardrails AI.
Probabilistic text-level validators; same category as Guardrails AI.

\section{Competitive Summary}

No independent evaluation comparing Pramanix against any of the above systems
on a common benchmark has been performed.  The verifiable technical
distinction is that Pramanix provides Z3 formal arithmetic verification with
per-invariant attribution and an Ed25519-signed audit chain --- capabilities
not present in any of the tools surveyed based on their public documentation.
Whether that constitutes a competitive advantage in any specific deployment
depends on workload characteristics.

A realistic production deployment in regulated industries may use these tools
in combination: OPA for authorization, Guardrails AI or LLM Guard for content
safety, and Pramanix for structured action verification.  These address
different problem classes and are not mutually exclusive.

%% ============================================================================
\chapter{Economics of Running}
%% ============================================================================

\section{Compute Cost Per Decision}

At 81.3~RPS single-threaded (from \fname{1m\_audit\_summary.json}), one CPU
core handles approximately 81 verifications per second in \code{sync} mode.
Each decision costs $\sim$12.3\,ms of wall-clock time including all overhead.

At the 18-worker architecture (from \fname{benchmarks/100m\_orchestrator\_fast.py}):
\begin{align*}
  18 \text{ workers} \times 80\ \text{RPS}  &= 1{,}440\ \text{aggregate RPS (minimum)} \\
  18 \text{ workers} \times 120\ \text{RPS} &= 2{,}160\ \text{aggregate RPS (maximum)}
\end{align*}

Projecting the single-core throughput from the Windows~11 laptop benchmark
(Table~\ref{tab:stability}) to a 32-vCPU \code{c5.8xlarge} (AWS, approximately
\$1.36/hour as of early 2026): a rough first-order estimate of approximately
\$0.0034 per million decisions in compute only.  This is a projection, not a
measurement on that instance type; actual throughput and cost on server
hardware require re-benchmarking on the target platform.

\section{Memory Budget}

Per worker: $\sim$50--80\,MiB peak Z3 heap $+$ 5\,MiB payload cache $+$
baseline Python process ($\sim$57\,MiB baseline RSS from the 1M benchmark).

$18 \text{ workers} \times 80\,\text{MiB peak} = \sim 1.4\,\text{GB RAM}$
for the full 100M/domain audit configuration. A production deployment with
4--8 workers requires 320--640\,MB, well within any modern server instance.

\section{SaaS Unit Economics}

\textbf{License precondition.}  The following analysis assumes a license
change from AGPL-3.0-only to a commercial or permissive license.  Under
AGPL-3.0~\cfrsec{13}, any organisation providing Pramanix as part of a
network service must make the complete corresponding source code of that
service available to its users.  A proprietary SaaS business model on this
codebase is not viable without a license change.  This is the primary
adoption blocker documented in \fname{reality.md} \cfrsec{4.2}.

Assuming a license change: at \$0.01 per 1,000 decisions (\$10/million),
gross margin relative to measured infrastructure cost (\$0.0034/million on a
single \code{c5.8xlarge} core in \code{sync} mode) is high.  The pricing
ceiling is set by the cost of a single incorrect decision in the target
domain.  These are rough projections from the single-machine 1M benchmark
(\fname{benchmarks/results/1m\_audit\_summary.json}).  Production economics
depend on operational overhead, multi-tenant isolation costs, policy
complexity, and invariant count not captured in that benchmark.

%% ============================================================================
\chapter{Security Threat Model}
%% ============================================================================

\section{The Eleven-Layer Defence Stack}

The complete stack, verified from source (\fname{reality.md} \cfrsec{2.1}):

\begin{table}[h]
\centering\small
\begin{tabular}{@{} c l p{7.5cm} @{}}
\toprule
\textbf{Layer} & \textbf{Location} & \textbf{Defence} \\
\midrule
1  & \fname{\_platform.py}                   & Alpine/musl rejection at import time \\
2  & \fname{translator/injection\_filter.py} & Pre-compiled regex, $\sim$25 injection patterns \\
3  & \fname{translator/\_sanitise.py}         & NFKC normalisation $+$ length limits \\
4  & \fname{validator.py}                    & Pydantic~v2 strict mode, zero implicit coercions \\
5  & \fname{translator/redundant.py}         & Dual-model LLM consensus \\
6  & \fname{guard\_pipeline.py}              & Post-consensus semantic check \\
7  & \fname{fast\_path.py}                   & Configurable $O(1)$ pre-screener, can only BLOCK \\
8  & \fname{solver.py}                       & Z3 with \code{rlimit} DoS guard \\
9  & \fname{translator/\_sanitise.py}         & Post-consensus injection heuristics \\
10 & \fname{guard.py}                        & Governance gates (Privilege, Oversight, IFC) \\
11 & \fname{crypto.py}                       & Ed25519 decision signing \\
\bottomrule
\end{tabular}
\caption{Eleven-layer defence stack (source: \fname{reality.md} \cfrsec{2.1})}
\end{table}

\section{Attack Vector Analysis}

\textbf{Prompt Injection}: The five-layer translator defence handles this. NFKC
normalisation collapses homoglyphs, regex patterns detect common overrides,
dual-model consensus requires both LLMs to agree. The Z3 layer is unreachable
via text injection because the policy is compiled Python DSL ---
\code{"ignore previous instructions"} cannot affect the Z3 formula.

\textbf{Malicious Z3 Constraints}: Both wall-clock timeout and \code{rlimit}
(elementary operation cap) are applied. \code{rlimit} prevents non-linear
expression DoS that can defeat wall-clock timeout via adversarial formula
construction.

\textbf{Worker Process Exploitation}: In \code{async-process} mode, the
HMAC-sealed IPC means a compromised worker cannot forge \code{allowed=True}
in the host process without the ephemeral key, which is stored only in the
host process and is not serializable.

\textbf{HMAC Key Compromise} (\fname{reality.md} \cfrsec{2.1}):
\fname{injection\_scorer.py}:305 uses \code{pickle.loads()} after HMAC
verification. If \code{PRAMANIX\_SCORER\_KEY} is compromised, this is an RCE
vector. HMAC key hygiene is a deployment requirement.

\textbf{JWT Algorithm Confusion} (fixed, \fname{CHANGELOG.md} \cfrsec{BUG-10}):
The CVE-2015-9235 family attack --- crafting a token with \code{alg: "RS256"}
and an HMAC-SHA256 signature --- was present and is now fixed. The header is
decoded and \code{alg} validated as \code{"HS256"} \emph{before} signature
computation.

\textbf{Timing Oracle} (fixed, \fname{CHANGELOG.md} \cfrsec{H-02}): ALLOW
responses returning faster than BLOCK responses allow binary-search probing to
identify violated constraints even with \code{redact\_violations=True}. Fixed:
timing pad applied unconditionally to every response before the ALLOW/BLOCK
branch.

\section{The Audit Non-Repudiation Chain}

Four layers of non-repudiation:

\begin{enumerate}
  \item \textbf{Decision hash}: SHA-256 over \code{decision\_id + status +
        allowed + violated\_invariants + explanation}. Deterministic. Any
        tampering is detectable.
  \item \textbf{Ed25519 signature}: Signs the decision hash. Offline
        verifiable with only the public key.
  \item \textbf{Merkle proof}: Proves inclusion in an unaltered batch without
        replaying all decisions.
  \item \textbf{ProvenanceChain}: HMAC-signed chain-of-custody binding each
        decision to policy version, model version, and IFC input labels.
\end{enumerate}

These four layers make it computationally infeasible to forge or alter a
decision record without either breaking Ed25519 or invalidating the Merkle
tree, \emph{given that the Ed25519 signing key has not been compromised}.  If
the private key is extracted from the host process or its key store, an
attacker with that key can produce valid signatures on fabricated decisions.
Key generation, storage, access control, and rotation for the Ed25519 private
key are the responsibility of the deploying organisation.  Pramanix provides
the signing and offline verification infrastructure; it does not manage key
lifecycle.

%% ============================================================================
\chapter{Academic and Research Grounding}
%% ============================================================================

\section{SMT Solving --- The Theoretical Foundation}

Satisfiability Modulo Theories (SMT) is the decision problem for logical
formulas with respect to background theories. The theories Pramanix uses:

\textbf{Linear Real Arithmetic (LRA)}: For financial constraints of the form
$\mathit{balance} - \mathit{amount} \geq \mathit{min\_reserve}$.  LRA is
decidable by quantifier elimination over real closed fields (Tarski, 1951).
Presburger arithmetic is a distinct decidable theory covering linear
arithmetic over non-negative integers; the two theories are different in
domain and in complexity.  Z3 uses the Simplex method for LRA satisfiability
checking, with polynomial average-case complexity for the linear constraint
patterns common in financial policy.

\textbf{Linear Arithmetic over Integers (LIA)}: For count-based constraints.
Also decidable.

\textbf{Boolean combinations via DPLL(T)}: Z3's architecture combines SAT
solving (propositional Boolean reasoning) with theory solvers. The DPLL SAT
solver handles the Boolean structure; the theory solver (Simplex for LRA)
handles the arithmetic atoms. For a constraint such as:
\begin{equation}
  (balance - amount \geq 0) \;\wedge\; (amount \leq daily\_limit) \;\wedge\;
  \neg frozen
  \label{eq:banking-policy}
\end{equation}
the Boolean skeleton $A \wedge B \wedge C$ is handled by DPLL; the arithmetic
atoms $A$ and $B$ are dispatched to the LRA solver; the Boolean atom $C$ is
handled natively.

\section{Formal Methods in Safety-Critical Systems}

Formal verification in safety-critical domains (avionics, nuclear, medical
devices) has used model checkers and theorem provers for decades. The NASA
Ames Formal Methods group, the Cambridge L4.verified project, and seL4
demonstrate that formal verification at scale is achievable in production
systems.

Pramanix applies this philosophy to AI agent actions. The argument is exactly
the same as in avionics: when the cost of an incorrect action is high enough,
probabilistic correctness is insufficient. A bank transfer that overdrafts an
account is not a statistic; it is an individual financial loss.

\section{Neuro-Symbolic AI}

The neuro-symbolic framing reflects the architecture: a neural component (LLM
for intent extraction) combined with a symbolic component (Z3 for formal
verification). This is a recognized research paradigm, with work from IBM
Research, MIT, and DeepMind on neuro-symbolic integration.

Pramanix's contribution in this implementation is a concrete, deployable
instantiation: the LLM extracts structured intent from natural language; the
symbolic verifier checks that intent against formal constraints using Z3. The dual-model consensus is a
neuro-symbolic defense: two independent neural extractors must agree before
the symbolic verifier is invoked. Disagreement itself is a signal --- ambiguous
or adversarial input produces consistent disagreement between models.

\section{The Z3 Literature}

Microsoft Research's Z3 has been continuously developed since 2007 (de~Moura
and Bj\o{}rner, \textit{Z3: An Efficient SMT Solver}, TACAS~2008). Key
properties exploited by Pramanix:

\begin{itemize}
  \item \textbf{DPLL(T) completeness}: For decidable theories (LRA, LIA), Z3
        either returns SAT with a model or UNSAT with a proof. No false
        positives.
  \item \textbf{Unsat cores}: \code{assert\_and\_track + unsat\_core()} extracts
        the minimal set of constraints responsible for unsatisfiability.
        Pramanix uses one constraint per solver to get complete attribution
        without relying on Z3's heuristic minimal-core extraction.
  \item \textbf{Python API / GIL release}: \code{z3-solver} is the official
        Python binding. The GIL is released during solver execution
        (\fname{docs/DECISIONS.md} \cfrsec{1}), enabling genuine thread-level
        parallelism.
\end{itemize}

%% ============================================================================
\chapter{Failure Log}
%% ============================================================================

\section{Architectures That Were Tried and Rejected}

\begin{itemize}
  \item \textbf{Recursive Merkle root construction}: The initial
        \code{MerkleAnchor.\_build\_root} was recursive. It hit Python's
        default call stack limit at $\sim$1,000 decisions.
        (\fname{docs/DECISIONS.md} \cfrsec{18})

  \item \textbf{\code{fork()} instead of \code{spawn()} for worker processes}:
        Fork copies parent file descriptors, memory-mapped files, and OS
        resources. Z3 has internal state --- forking after Z3 initialization
        produced undefined behavior. (\fname{docs/DECISIONS.md} \cfrsec{4})

  \item \textbf{\code{threading.local} for resolver cache}: Thread-scoped.
        Under asyncio, Task~B could see Task~A's resolved field values --- a
        data bleed between users. (\fname{docs/DECISIONS.md} \cfrsec{6})

  \item \textbf{Single solver for unsat attribution}: Z3's \code{unsat\_core()}
        on a multi-assertion solver does not always return the minimal core.
        One solver per invariant guarantees complete, deterministic attribution.
        (\fname{docs/DECISIONS.md} \cfrsec{3})

  \item \textbf{Eval/exec in the transpiler}: String-based formula generation
        is an RCE surface. Rejected entirely. (\fname{docs/DECISIONS.md}
        \cfrsec{9})

  \item \textbf{Floating-point direct to Z3}: IEEE~754 rounding errors.
        \code{Decimal(str(v)).as\_integer\_ratio()} was required.
        (\fname{docs/DECISIONS.md} \cfrsec{10})
\end{itemize}

\section{Security Bugs Found During Development}

\begin{itemize}
  \item \textbf{Silent stub exploitation (H-03, H-04)}: LangChain and CrewAI
        integrations both had \code{execute\_fn=None} defaulting to
        \code{lambda i: "OK"}. Every guarded action returned success without
        executing real logic. Fixed by raising \code{NotImplementedError}.

  \item \textbf{Async fail-safe bypass (C-01)}: \code{verify\_async()} had
        \code{except Exception: pass} in the size-check path. Unserializable
        payloads silently bypassed the size gate. Any exception in the size
        check must produce \code{Decision.error(allowed=False)}.

  \item \textbf{Timing oracle (H-02)}: ALLOW and BLOCK responses having
        different latencies is a side-channel attack. The timing pad was applied
        only to BLOCK. Fixed: timing pad applied unconditionally before the
        ALLOW/BLOCK branch.

  \item \textbf{JWT algorithm confusion (BUG-10)}: CVE-2015-9235 family.
        Decoding the payload before validating the algorithm allows crafting
        tokens with mismatched algorithms. Fixed.
\end{itemize}

\section{The Hardest Bugs}

\textbf{Z3 context threading}: Z3 \code{Context} objects are thread-local C
objects. They cannot be shared between threads. The GIL is released during
solving --- meaning Z3 operations in multiple threads run truly concurrently
in C code. Using shared contexts between threads produces undefined behavior.
Solution: thread-local \code{z3.Context()} created once per OS thread, never
destroyed. The adversarial test (\fname{test\_z3\_context\_isolation.py}) uses
\code{threading.Barrier} to force 10 concurrent simultaneous solver executions.

\textbf{The \code{ConstraintExpr.\_\_bool\_\_} trap}: \code{expr1 and expr2}
in Python short-circuits at the Python level. If \code{\_\_bool\_\_} returned
\code{True}, the \code{and} expression would return \code{expr2}, silently
discarding \code{expr1}. The fix --- raising \code{PolicyCompilationError} ---
means the common mistake (using Python \code{and}/\code{or} instead of
\code{\&}/\code{|}) is immediately detected at policy compile time.

%% ============================================================================
\chapter{Real World Simulation Results}
%% ============================================================================

\section{The 1M Decision Stability Test}

The most important stability result is in
\fname{benchmarks/results/1m\_audit\_summary.json}. One million decisions at
81.3 average~RPS, with memory and GC tracked throughout.

Key finding: Z3 heap behavior is oscillatory, not leaking. The
\fname{1m\_audit\_full.log} shows RSS spiking between 57\,MiB and 74\,MiB
throughout the entire run, then settling. The net growth over 1M decisions:
2.805\,MiB. GC gen0 collected 6 additional times; gen1 and gen2 unchanged.
This is bounded, stable behavior.

The P99.9 of 153.848\,ms and P99.99 of 270.578\,ms indicate rare but real
tail latency events. The max of 1,565.746\,ms is a single outlier --- likely
GC pause or OS scheduling jitter. For production deployments, a
\code{min\_response\_ms} padding budget and circuit breaker should handle these
tail events.

\section{Finance Domain Stress Test}

\fname{benchmarks/100m\_domain\_policies.py} documents the Finance domain
benchmark policy (\code{Finance100MPolicy}) targeting:
\begin{itemize}
  \item 20\,\% overdraft (BLOCK on \code{non\_negative\_balance})
  \item 10\,\% high risk (BLOCK on \code{risk\_score\_below\_threshold})
  \item $\sim$5\,\% overlap (double BLOCK)
  \item Target $\sim$25\,\% BLOCK rate overall
\end{itemize}

This mix exercises both Phase~1 (SAT, $\sim$75\,\% of decisions) and Phase~2
(UNSAT attribution, $\sim$25\,\% of decisions). A policy tuned only for
common-case SAT performance that degrades badly on UNSAT would not survive
this benchmark.

\section{Memory Growth Model}

The benchmark architecture's \code{max\_decisions\_per\_worker = 10{,}000}
Guard recycling strategy was empirically derived. The Guard recycle strategy:

\begin{enumerate}
  \item After 10,000 decisions, the old Guard is replaced with a new one
  \item The old Guard's worker pool is drained (background thread,
        \code{\_drain\_executor})
  \item Z3 heap of the old Guard is freed by Python's GC when the Guard
        object's reference count drops to zero
\end{enumerate}

This keeps Z3 heap bounded at ${<}$50\,MiB growth per worker.

%% ============================================================================
\chapter{Dependency Risk Analysis}
%% ============================================================================

\section{Z3 API Change Risk}

The constraint in \fname{pyproject.toml}: \code{z3-solver = "\^{}4.12"}.
The \code{\^{}} operator means \code{>=4.12.0, <5.0.0}. Z3's 4.x line has
been stable since 2018. The parts Pramanix uses --- \code{.check()},
\code{.model()}, \code{.unsat\_core()}, \code{z3.RatVal()},
\code{z3.Context()} --- are stable public API.

Mitigation: the transpiler is the only file calling Z3 directly. A Z3 API
change requires updating only \fname{transpiler.py} and \fname{solver.py}.
The rest of the codebase is Z3-agnostic.

\section{Z3 Swap Difficulty}

A swap to another SMT solver (e.g., cvc5, MathSAT) would require:
\begin{enumerate}
  \item A new transpiler file emitting the target solver's AST
  \item A new solver module wrapping the target solver's Python API
  \item Compatibility verification for \code{RatVal} exact rationals,
        \code{assert\_and\_track + unsat\_core()} attribution, and GIL release
        during solve
\end{enumerate}

The DSL and all user-facing policy code would be unchanged. Estimated effort:
500--800 lines of code in transpiler and solver, plus thorough re-testing of
the exact rational arithmetic behavior.

\section{LLM Deprecation Risk}

The translator layer supports seven providers. If any single provider
deprecates its API, only one adapter file needs updating. The
\code{OpenAICompatTranslator} covers OpenAI, Azure OpenAI, and any
OpenAI-compatible endpoint, reducing single-provider lock-in.

\section{Single Points of Failure}

The only genuine single point of failure is Z3. If Z3 is unavailable, the
entire verification stack fails. The musl check (\fname{\_platform.py}) fails
fast at import time. There is no Z3 fallback --- by design, the system is
fail-closed when Z3 is unavailable.

%% ============================================================================
\chapter{Regulatory Mapping}
%% ============================================================================

\textbf{Disclaimer.}  This chapter describes how Pramanix SDK components
relate to specific regulatory text.  These are engineering observations by
the author, not legal analysis.  Regulatory compliance is a legal and
operational determination that depends on jurisdiction, deployment context,
business processes, and regulatory examiner interpretation.  Nothing in this
chapter constitutes legal advice, and use of the Pramanix SDK alone does not
confer regulatory compliance on any deployment.  Organisations subject to
BSA/AML, HIPAA, SOX, EU AI Act, or other regulations must obtain qualified
legal counsel before making compliance claims.

\section{Financial Regulation}

The \code{ComplianceReporter} in \fname{helpers/compliance.py} maps
\code{violated\_invariants} Z3 labels to structured regulatory citations:

\begin{table}[h]
\centering\small
\begin{tabular}{@{} l l l @{}}
\toprule
\textbf{Domain} & \textbf{Regulation} & \textbf{Reference} \\
\midrule
BSA/AML   & Bank Secrecy Act, Anti-Structuring & 31~CFR~\S\,1020.320(a)(2) \\
OFAC/SDN  & Sanctions screening                & 50~CFR~\S\,598 \\
SEC       & Wash Sale                          & IRC~\S\,1091 \\
SOX       & Financial controls reporting       & 15~U.S.C.~\S\,7241 \\
Basel~III & Capital adequacy                   & BCBS~189 \\
\bottomrule
\end{tabular}
\caption{Financial regulatory mappings in \code{ComplianceReporter}
  (\fname{helpers/compliance.py})}
\end{table}

\section{Healthcare Regulation (HIPAA)}

\code{PHILeastPrivilege}, \code{BreakGlassAuth}, \code{ConsentActive},
\code{DosageGradientCheck}, \code{PediatricDoseBound} in
\fname{primitives/healthcare.py} map to HIPAA 45~CFR~\S\,164.

\textbf{Audit trail for HIPAA}: Every PHI access decision produces a
\code{Decision} with \code{decision\_hash}, Ed25519 signature, and Merkle
inclusion proof.  These tamper-evident signed records can be incorporated
into a HIPAA~\S\,164.312(b) audit trail programme.  They do not in isolation
satisfy HIPAA audit logging requirements, which also encompass user
login/logout events, system activity logs from all components handling ePHI,
and data access and modification records across the full application stack.
Pramanix covers the AI agent action layer of an audit trail, not the complete
HIPAA audit logging obligation.

\textbf{Archive encryption (resolved)}: \code{EncryptedArchiveWriter}
(Phase~10) provides AES-256-GCM at-rest encryption for all Merkle archive
segments.  PHI-adjacent audit logs are no longer stored in plaintext.
Deployments pass \code{archive\_writer=EncryptedArchiveWriter(key)} at
\code{MerkleArchiver} construction.  HIPAA \S164.312(a)(2)(iv) (encryption
and decryption) is addressable at the SDK layer.

\section{EU AI Act}

The EU AI Act (Regulation (EU) 2024/1689) classifies high-risk AI systems
(Annex~III, including financial services AI, healthcare AI, and critical
infrastructure AI) with requirements including Article~12 record-keeping.
Pramanix's signed decisions, Merkle anchoring, and ProvenanceChain directly
support Article~12. The \code{ComplianceReporter} could be extended to
produce EU AI Act Article~13 transparency reports.

\section{SOC~2 Type~II}

\begin{itemize}
  \item \textbf{CC6.1 (Logical access)}: JWT identity boundary with
        HMAC-SHA256 verification and \code{sub} claim validation
  \item \textbf{CC7.2 (System operations monitoring)}: 6 configurable audit
        sinks (Kafka, S3, Splunk, Datadog)
  \item \textbf{CC7.3 (Incident detection)}: \code{AdaptiveCircuitBreaker}
        state machine with \code{ISOLATED} state requiring manual reset
  \item \textbf{Confidentiality (C1)}: \code{\_RedactSecretsProcessor} in
        structlog pipeline, IFC labels
  \item \textbf{CC6.6 (Encryption at rest)}: \textit{[Resolved — Phase~10]}
        \code{EncryptedArchiveWriter} provides AES-256-GCM encryption for all
        Merkle archive segments.  Pass the writer at construction;
        no external tool chain required.
\end{itemize}

%% ============================================================================
\chapter{Scaling Story}
%% ============================================================================

\section{10$\times$ Scale: Current Architecture}

At 10$\times$ current load (810~RPS), a 10-worker configuration:
\begin{itemize}
  \item $10 \times 80\text{--}120\ \text{RPS} = 800\text{--}1{,}200$
        aggregate~RPS
  \item RAM: $10 \times 80\,\text{MiB} = 800\,\text{MiB}$
  \item CPU: 10 dedicated cores (or hardware threads)
\end{itemize}

No architectural changes required. The circuit breaker begins shedding load
when P99 latency exceeds \code{shed\_latency\_threshold\_ms}.

\section{100$\times$ Scale: Horizontal Z3 Worker Pools}

At 100$\times$ (8,100~RPS), a single machine with 18 workers approaches its
ceiling ($\sim$2,160~RPS). Scaling to 8,100~RPS requires 4--6 machines or a
Kubernetes deployment.

Architecture:
\begin{itemize}
  \item K8s \code{PramanixMiddleware} admission webhook $\to$ shared policy
        hash check at startup
  \item StatelessGuard instances per pod (no inter-pod state, policies are
        immutable)
  \item Redis \code{DistributedCircuitBreaker} for coordinated load shedding
        across pods
  \item Redis \code{RedisExecutionTokenVerifier} for cross-pod single-use
        token enforcement
\end{itemize}

The policy hash check (\code{GuardConfig.expected\_policy\_hash}) is the
critical invariant: all replicas must run the same policy version. A SHA-256
mismatch at construction time raises \code{ConfigurationError} --- prevents
split-brain policy enforcement during rolling deploys.

\section{1000$\times$ Scale: What Breaks First}

At 1,000$\times$ (81,000~RPS), several architectural limits are hit:

\begin{itemize}
  \item \textbf{Z3 memory}: $\sim$80\,MiB peak per worker $\times$ 750 workers
        $= \sim$60\,GB peak RAM. Horizontal pod scaling with smaller per-pod
        worker counts mitigates.

  \item \textbf{Redis circuit breaker TOCTOU} --- \textit{resolved}.
        \code{RedisDistributedBackend.set\_state} uses \code{WATCH} +
        \code{MULTI}/\code{EXEC} optimistic locking with
        \code{\_MERGE\_MAX\_RETRIES=3} retries, eliminating the race
        without Lua scripting.

  \item \textbf{Audit sink throughput}: Kafka sink's 10,000-message queue
        overflows if backend throughput drops below 81,000 messages/second.
        Fix: increase queue size or add Kafka partitioning.

  \item \textbf{Execution token consumed-set}: In-memory per-process set.
        \code{RedisExecutionTokenVerifier} becomes mandatory --- and Redis
        becomes a SPOF.
\end{itemize}

%% ============================================================================
\chapter{Developer Experience Decisions}
%% ============================================================================

\section{Python vs.\ Rust/Go}

The decision to use Python is implicit in the architecture. AI agent frameworks
(LangChain, AutoGen, CrewAI) are Python-first. The SDK must integrate without
a language boundary. \code{z3-solver} is the official Python binding --- Z3's
Rust bindings (\code{z3} crate) are less mature and do not have feature parity.
The policy DSL is natural Python: \code{E(balance) - E(amount) >= 0}. \code{mypy}
strict mode with \code{disallow\_untyped\_defs = true} provides close to
Rust-level type safety. Z3 releases the Python GIL during solve, enabling
genuine parallel solving in thread mode without Python-level synchronization.

\section{Poetry vs.\ pip}

\fname{pyproject.toml} uses Poetry's build backend
(\code{poetry-core>=1.7.0}). The lockfile-based dependency resolution prevents
the transitive dependency conflicts that bare pip installs produce with 31
optional extras spanning AI, cloud, database, cryptography, and compliance
ecosystems.

\section{AGPL-3.0 vs.\ Apache~2.0 vs.\ MIT}

AGPL-3.0 is documented in \fname{pyproject.toml}:8 and analyzed in
\fname{reality.md} \cfrsec{4.2} as the first absolute blocker:

\begin{quote}
AGPL-3.0 License: No enterprise legal team approves an AGPL Python
dependency. Blocks fintech, healthcare, and enterprise B2B adoption.
\end{quote}

AGPL's copyleft requirement: any software that uses an AGPL library and is
provided as a network service must release its source code. For a bank or
hospital using Pramanix in production, this means releasing their entire
application source. Most enterprise legal teams in regulated industries do
not accept this condition without an explicit exception review and relicensing
agreement.

Apache~2.0 is the correct choice: patent protection (important for Z3 usage),
enterprise-approved legal language, and allows proprietary use without source
release.

\section{DSL Syntax Alternatives}

The DSL design (\fname{expressions.py}) chose Python operator overloading
(\code{\&}, \code{|}, \code{\textasciitilde}, \code{>=}, \code{<=},
\code{>}, \code{<}, \code{==}, \code{!=}) over alternatives:

\begin{itemize}
  \item \textbf{String-based DSL}: Rejected --- requires \code{eval}/\code{parse}
        (injection vector), no type checking, no IDE support.
  \item \textbf{YAML-based policy}: Rejected --- cannot express Python-typed
        cross-field relationships without a custom schema language.
  \item \textbf{Rego/Datalog-style}: Rejected --- lacks Z3's arithmetic, would
        require embedded parser.
  \item \textbf{JSON Schema-based}: Rejected --- JSON Schema is not a constraint
        specification language; it validates types, not cross-field relationships.
\end{itemize}

The Python operator overloading approach provides IDE support, type checking,
and Python-native ergonomics. The \code{\_\_bool\_\_} raising provides
immediate feedback when the common mistake (\code{and} instead of \code{\&})
is made at policy compile time.

%% ============================================================================
\chapter{Personal Reflection}
%% ============================================================================

\section{What This Project Actually Is}

Pramanix is an attempt to answer one specific engineering question: can
production-quality formal verification be made accessible to Python developers
who build AI systems?

The answer, based on the evidence of this codebase, is yes --- with significant
caveats. The caveats are not embarrassments; they are the research findings.

\section{What Required the Most Learning}

\textbf{Z3 internals}: Thread-local contexts, GIL release semantics,
\code{assert\_and\_track} vs.\ \code{add()}, \code{RatVal} vs.\ \code{IntVal}
vs.\ \code{RealVal}, \code{rlimit} as a DoS defense. None of this is in the
Z3 Python tutorial. It required reading Z3's C++ source code, the DPLL(T)
algorithm papers, and extensive empirical testing under adversarial conditions.

\textbf{The memory model}: Understanding why Z3 heap oscillates rather than
leaks, how Python's GC interacts with C++ allocated objects via
\code{z3-solver}, when to recycle a Guard instance --- these required building
the 1M-decision benchmark and reading the log output at the line level.

\textbf{Exact arithmetic in SMT}: The \code{Decimal.as\_integer\_ratio()}
insight was not obvious. IEEE~754 floating-point passed to Z3's \code{z3.Q(...)}
produces floating-point arithmetic in Z3, not real arithmetic.
\code{z3.RatVal(numerator, denominator)} produces exact rational arithmetic.
This distinction matters for financial invariants and took significant
experimentation to understand and verify.

\textbf{Security reasoning in depth}: The JWT algorithm confusion bug
(BUG-10), the async fail-safe bypass (C-01), the timing oracle (H-02), the
worker HMAC seal --- these are not bugs that surface in happy-path testing.
They surface in adversarial analysis. Building the threat model before the
implementation would have caught several earlier.

\section{What the Audit Taught}

The \fname{reality.md} document is the most important artifact in this
repository. It is an honest, line-by-line audit of every claim made about the
system against the actual source code. Several claimed features were initially
described incorrectly; the audit corrected them. Several real gaps were found
that the initial implementation missed.

The practice of writing this kind of audit --- before any external release ---
is what separates engineering from marketing. The gaps documented in
\fname{reality.md} are not weaknesses to hide; they are the work remaining
to be done.

Three truths from the audit that summarize the project:

\begin{enumerate}
  \item \textbf{The technical differentiation is verifiable from source.}
        Eleven-layer defence stack, 38 pre-built policy primitives, compliance
        reporter mapping Z3 violation labels to regulatory citations in six
        frameworks, audit chain with Ed25519 signing and Merkle anchoring.
        Every component cited here has a specific source file reference.  The
        self-audit in \fname{reality.md} verified these against actual code.

  \item \textbf{The license is the primary blocker.}  Every claim about
        enterprise deployment potential is gated on AGPL-3.0.  Switch to
        Apache~2.0 or a commercial license before any external release.

  \item \textbf{The ComplianceReporter and ShadowEvaluator need documentation
        prominence.}  A compliance reporter that maps a Z3 violation label to
        a BSA/AML CFR citation with PDF export.  A shadow evaluator that runs
        candidate policy non-blocking against live traffic before promotion.
        These capabilities are not obvious from the README.  They need
        explicit front-page documentation.
\end{enumerate}

\section{What Comes Next}

The fail-closed contract must never be relaxed. The Z3 path to
\code{allowed=True} must remain the only path.

The license must change. The Redis circuit breaker race must be fixed. The
logging split-brain must be unified.  \textit{[Done --- Phase~10]} The Merkle
archive encryption is now pluggable (\code{EncryptedArchiveWriter}).

After those three remaining changes (license, Redis race, logging split-brain),
Pramanix is ready for external release. Not because
the system is perfect --- the 14 known gaps in \fname{docs/KNOWN\_GAPS.md}
document what remains. But because every gap is known, named, and honest, and
the core contract has not been broken.

That is what production readiness means: not the absence of gaps, but the
honest accounting of them.

%% ============================================================================
%%  APPENDICES
%% ============================================================================

\appendix

\chapter{File Reference Map}
\label{app:references}

\begin{longtable}{@{} p{5.5cm} p{5cm} p{3.5cm} @{}}
\toprule
\textbf{Claim} & \textbf{File} & \textbf{Location} \\
\midrule
\endfirsthead
\multicolumn{3}{l}{\footnotesize\textit{(continued from previous page)}} \\
\toprule
\textbf{Claim} & \textbf{File} & \textbf{Location} \\
\midrule
\endhead
\bottomrule
\multicolumn{3}{r}{\footnotesize\textit{(continued on next page)}} \\
\endfoot
\bottomrule
\endlastfoot
Fail-closed contract        & \fname{docs/DECISIONS.md}     & \cfrsec{2} \\
Two-phase solver            & \fname{docs/DECISIONS.md}     & \cfrsec{3} \\
Exact arithmetic            & \fname{docs/DECISIONS.md}     & \cfrsec{10} \\
HMAC-sealed IPC             & \fname{docs/DECISIONS.md}     & \cfrsec{15} \\
ContextVar isolation        & \fname{docs/DECISIONS.md}     & \cfrsec{6} \\
Alpine/musl rejection       & \fname{docs/DECISIONS.md}     & \cfrsec{12} \\
Iterative Merkle            & \fname{docs/DECISIONS.md}     & \cfrsec{18} \\
Timing pad design           & \fname{docs/DECISIONS.md}     & \cfrsec{19} \\
P50/P95/P99 latency         & \fname{benchmarks/results/latency\_results.json} & root \\
1M stability results        & \fname{benchmarks/results/1m\_audit\_summary.json} & root \\
100M architecture           & \fname{benchmarks/100m\_orchestrator\_fast.py} & header \\
38 primitives               & \fname{reality.md}            & \cfrsec{2.2} \\
AGPL blocker                & \fname{reality.md}            & \cfrsec{4.2} \\
Redis TOCTOU race           & \fname{reality.md}            & \cfrsec{4.1} \\
Logging split-brain         & \fname{reality.md}            & \cfrsec{4.2} \\
4,002 tests (post-upgrade)  & \fname{reality.md}            & \cfrsec{3.1} \\
3,916 passing (Phase~10)    & \fname{tests/unit/test\_archiver\_coverage\_gaps.py} & full file \\
3,916 passing (Phase~11)    & \fname{tests/unit/test\_interceptors\_real.py} & \code{TestKafkaConsumerRealPaths} \\
EncryptedArchiveWriter      & \fname{src/pramanix/audit/archiver.py} & class EncryptedArchiveWriter \\
amount\_field parameter     & \fname{src/pramanix/translator/\_sanitise.py} & injection\_confidence\_score \\
structlog isolation fix     & \fname{tests/unit/test\_guard\_dark\_paths.py} & \code{\_capture\_log\_output} \\
Stability tiers             & \fname{docs/PUBLIC\_API.md}   & Stability Tiers \\
9 integrations              & \fname{reality.md}            & \cfrsec{2.4} \\
ComplianceReporter          & \fname{reality.md}            & \cfrsec{2.4} \\
ShadowEvaluator             & \fname{reality.md}            & \cfrsec{2.4} \\
BUG-10 JWT fix              & \fname{docs/CHANGELOG.md}     & \cfrsec{BUG-10} \\
H-02 timing oracle fix      & \fname{docs/CHANGELOG.md}     & \cfrsec{H-02} \\
C-01 async bypass fix       & \fname{docs/CHANGELOG.md}     & \cfrsec{C-01} \\
H-03 LangChain stub         & \fname{docs/CHANGELOG.md}     & \cfrsec{H-03} \\
H-04 CrewAI stub            & \fname{docs/CHANGELOG.md}     & \cfrsec{H-04} \\
BUG-03 OversightRecord      & \fname{docs/CHANGELOG.md}     & \cfrsec{BUG-03} \\
\end{longtable}

\chapter{Dependency Inventory}
\label{app:deps}

\textbf{Core} (always installed with \code{pip install pramanix}):

\begin{table}[h]
\centering
\begin{tabular}{@{} l l p{7cm} @{}}
\toprule
\textbf{Package} & \textbf{Version} & \textbf{Purpose} \\
\midrule
\code{pydantic}  & \code{\textasciicircum{}2.5}  & Strict-mode input validation \\
\code{z3-solver} & \code{\textasciicircum{}4.12} & SMT formal verification engine \\
\code{structlog} & \code{\textasciicircum{}23.2} & Structured logging with secrets redaction \\
\bottomrule
\end{tabular}
\caption{Core dependencies (source: \fname{pyproject.toml})}
\end{table}

\textbf{Optional extras} (31 total, from \fname{pyproject.toml}):

\begin{itemize}
  \item \code{translator}: \code{httpx}, \code{openai}, \code{anthropic},
        \code{tenacity}
  \item \code{otel}: \code{opentelemetry-sdk},
        \code{opentelemetry-exporter-otlp-proto-grpc}
  \item \code{fastapi}: \code{fastapi}, \code{starlette}, \code{httpx}
  \item \code{crypto}: \code{cryptography} (Ed25519 signing)
  \item \code{aws} / \code{azure} / \code{gcp} / \code{vault}: cloud key
        providers
  \item \code{kafka}: \code{confluent-kafka} (audit sink)
  \item \code{metrics}: \code{prometheus-client}
  \item \code{performance}: \code{orjson}
  \item \code{sklearn}: \code{scikit-learn} (\code{CalibratedScorer})
  \item Plus 20 more for specific framework integrations
\end{itemize}

%% ============================================================================
%%  BIBLIOGRAPHY
%% ============================================================================

\chapter*{References}
\addcontentsline{toc}{chapter}{References}

\begin{enumerate}[label={[\arabic*]}]
  \item L.~de~Moura and N.~Bj\o{}rner. \textit{Z3: An Efficient SMT Solver}.
        In \textit{Tools and Algorithms for the Construction and Analysis of
        Systems (TACAS~2008)}, Lecture Notes in Computer Science, vol.~4963,
        pp.~337--340. Springer, Berlin, Heidelberg, 2008.

  \item D.~Kroening and O.~Strichman. \textit{Decision Procedures: An
        Algorithmic Point of View}. 2nd~ed. Springer, 2016.

  \item R.~Nieuwenhuis, A.~Oliveras, and C.~Tinelli. Solving SAT and SAT
        Modulo Theories: From an Abstract Davis--Putnam--Logemann--Loveland
        Procedure to DPLL(T). \textit{Journal of the ACM}, 53(6):937--977,
        2006.

  \item Open Policy Agent. \textit{Policy Language}. Open Policy Agent
        Documentation. \url{https://www.openpolicyagent.org/docs/latest/}.
        Accessed May 2026.

  \item National Institute of Standards and Technology. \textit{Artificial
        Intelligence Risk Management Framework (AI RMF 1.0)}. NIST AI 100-1.
        January 2023.

  \item European Parliament and the Council of the European Union. Regulation
        (EU) 2024/1689 laying down harmonised rules on artificial intelligence
        (Artificial Intelligence Act). \textit{Official Journal of the European
        Union}, L series, 2024.

  \item G.~Katz, C.~Barrett, D.~Dill, K.~Julian, and M.~Kochenderfer.
        Reluplex: An Efficient SMT Solver for Verifying Deep Neural Networks.
        In \textit{Computer Aided Verification (CAV~2017)}, Lecture Notes in
        Computer Science, vol.~10426. Springer, 2017.

  \item seL4 Foundation. \textit{The seL4 Microkernel: An Introduction}.
        \url{https://sel4.systems/Info/Docs/seL4-manual-latest.pdf}.
        Accessed May 2026.
\end{enumerate}

\vfill
\begin{center}
  \small\itshape
  Every claim in this document is grounded in source code, benchmark output,\\
  or documentation that exists in the repository as of the audit date 2026-05-12.\\
  Where code and documentation disagree, the code is the authority.
\end{center}

\end{document}
